% Source: Evans, "Partial Differential Equations" (2nd ed., AMS Graduate Studies in Mathematics vol. 19, 2010),
% Chapter 11 "Systems of Conservation Laws", PDF pages 569-614 of MathBooks/Evans Partial Differential Equations(1).pdf.
% Transcribed page-by-page by vision subagents, then merged and restructured into
% leanblueprint conventions (semantic \label, bracketed titles, \uses dependency edges) below.

\section{Introduction}\label{sec:conservation-laws-introduction}

In this final chapter we study \textit{systems} of nonlinear, divergence structure first-order hyperbolic PDE, which arise as models of \textit{conservation laws}.

\textbf{Physical interpretation.} In the most general circumstance we would like to investigate a vector function

$$\mathbf{u} = u(x,t) = (u^1(x,t), \ldots, u^m(x,t)) \quad (x \in \mathbb{R}^n, t \geq 0),$$

the components of which are the densities of various conserved quantities in some physical system under investigation. Given then any smooth, bounded region $U \subset \mathbb{R}^n$, we note that the integral

$$(1) \qquad \int_U \mathbf{u}(x,t) \, dx$$

represents the total amount of these quantities within $U$ at time $t$. Now conservation laws typically assert that the rate of change within $U$ is governed by a \textit{flux} function $\mathbf{F}: \mathbb{R}^m \to M^{m \times n}$, which controls the rate of loss or increase of $\mathbf{u}$ through $\partial U$. Otherwise stated, it is appropriate to assume for each time $t$

$$(2) \qquad \frac{d}{dt} \int_U \mathbf{u} \, dx = - \int_{\partial U} \mathbf{F}(\mathbf{u})\nu \, dS,$$

$\nu$ denoting as usual the outward unit normal along $U$. Rewriting (2), we deduce

$$(3) \qquad \int_U \mathbf{u}_t \, dx = - \int_{\partial U} \mathbf{F}(\mathbf{u})\nu \, dS = - \int_U \text{div} \mathbf{F}(\mathbf{u}) \, dx.$$

As the region $U \subset \mathbb{R}^n$ was arbitrary, we derive from (3) this initial-value problem for a \textit{general system of conservation laws}:

$$(4) \qquad \begin{cases} \mathbf{u}_t + \text{div} \mathbf{F}(\mathbf{u}) = 0 & \text{in } \mathbb{R}^n \times (0, \infty) \\ \mathbf{u} = \mathbf{g} & \text{on } \mathbb{R}^n \times \{t = 0\}, \end{cases}$$

the given function $\mathbf{g} = (g^1, \ldots, g^m)$ describing the initial distribution of $\mathbf{u} = (u^1, \ldots, u^m)$.

At present a good mathematical understanding of problem (4) is largely unavailable (but see Zheng \cite{Zheng2001}). For this reason we shall henceforth consider instead the initial-value problem for a \textit{system of conservation laws in one space dimension}:

$$(5) \qquad \begin{cases} \mathbf{u}_t + \mathbf{F}(\mathbf{u})_x = 0 & \text{in } \mathbb{R} \times (0, \infty) \\ \mathbf{u} = \mathbf{g} & \text{on } \mathbb{R} \times \{t = 0\}, \end{cases}$$

where $\mathbf{F}: \mathbb{R}^m \to \mathbb{R}^m$ and $\mathbf{g}: \mathbb{R} \to \mathbb{R}^m$ are given and $\mathbf{u}: \mathbb{R} \times [0, \infty) \to \mathbb{R}^m$ is the unknown, $\mathbf{u} = \mathbf{u}(x,t)$. We call $\mathbb{R}^m$ the \textit{state space}, and write

$$\mathbf{F} = \mathbf{F}(z) = (F^1(z), \ldots, F^m(z)) \quad (z \in \mathbb{R}^m)$$

for the smooth flux function.

We intend to study the solvability of problem (5), properties of its solutions, etc.

\begin{example}[The $p$-system]
\label{ex:p-system-definition}
\dcref{ch11:11.1-ex-1}
\group{Conservation Laws}
The $p$-system is this collection of two conservation laws:

$$(6) \qquad \begin{cases} u^1_t - u^2_x = 0 & \text{(compatibility condition)} \\ u^2_t - p(u^1)_x = 0 & \text{(Newton's law)}, \end{cases}$$

in $\mathbb{R}\times(0,\infty)$, where $p:\mathbb{R}\to\mathbb{R}$ is given. Here

$$(7) \qquad \mathbf{F}(z) = (-z_2, -p(z_1))$$

for $z=(z_1,z_2)$. The $p$-system arises as a rewritten form of the scalar \textit{nonlinear wave equation}

$$(8) \qquad u_{tt} - (p(u_x))_x = 0 \quad \text{in } \mathbb{R}\times(0,\infty).$$

Taking $u^1 := u_x$, $u^2 := u_t$, we obtain the system (6), with the stated interpretations. $\square$
\end{example}

\begin{example}[Euler's equations for compressible gas flow]
\label{ex:euler-equations-1d-system}
\dcref{ch11:11.1-ex-2}
\group{Conservation Laws}
\textit{Euler's equations} for compressible gas flow in one dimension are

$$(9) \qquad \begin{cases} \rho_t + (\rho v)_x = 0 & \text{(conservation of mass)} \\ (\rho v)_t + (\rho v^2+p)_x = 0 & \text{(conservation of momentum)} \\ (\rho E)_t + (\rho E v + pv)_x = 0 & \text{(conservation of energy)} \end{cases}$$

in $\mathbb{R}\times(0,\infty)$. Here $\rho$ is the \textit{mass density}, $v$ the \textit{velocity}, and $E$ the \textit{energy density} per unit mass. We assume

$$E = e + \frac{v^2}{2},$$

where $e$ is the \textit{internal energy} per unit mass and the term $\frac{v^2}{2}$ corresponds to the kinetic energy per unit mass. The letter $p$ in (9) denotes the \textit{pressure}. We assume $p$ is a known function

$$(10) \qquad p = p(\rho,e)$$

of $\rho$ and $e$; formula (10) is a \textit{constitutive relation}. Writing $\mathbf{u}=(u^1,u^2,u^3)=(\rho,\rho v,\rho E)$, we check (9) is a system of conservation laws of the requisite form
$$\mathbf{u}_t + \mathbf{F}(\mathbf{u})_x = 0 \quad \text{in } \mathbb{R}\times(0,\infty)$$
for $\mathbf{F}=(F^1,F^2,F^3)$,

$$(11) \qquad \begin{cases} F^1(z) = z_2 \\ F^2(z) = \frac{(z_2)^2}{z_1} + p\left(z_1, \frac{z_3}{z_1} - \frac{1}{2}\left(\frac{z_2}{z_1}\right)^2\right) \\ F^3(z) = \frac{z_2 z_3}{z_1} + p\left(z_1, \frac{z_3}{z_1} - \frac{1}{2}\left(\frac{z_2}{z_1}\right)^2\right)\frac{z_2}{z_1}, \end{cases}$$

where $z=(z_1,z_2,z_3)$, $z_1 > 0$. $\square$
\end{example}

\begin{example}[The shallow water equations]
\label{ex:shallow-water-equations}
\dcref{ch11:11.1-ex-3}
\group{Conservation Laws}
The one-dimensional \textit{shallow water equations} are
$$\begin{cases} \phi_t + (v\phi)_x = 0 & \text{(conservation of mass)} \\ v_t + (v^2/2+\phi)_x = 0 & \text{(conservation of momentum)} \end{cases}$$
in $\mathbb{R}\times[0,\infty)$. In these equations $v$ is the horizontal velocity, $\phi$ denotes $gh$, where $h$ is the height and $g$ the gravitational constant, and
$$\mathbf{F}(z) = \left(z_1z_2, \frac{(z_2)^2}{2} + z_1\right). \qquad \square$$
\end{example}

\subsection{Integral solutions}\label{sec:integral-solutions-systems}

The great difficulty in this subject is discovering a proper notion of weak solution for the initial-value problem (5). We have already encountered similar issues in \S3.4 in our study of the much simpler case of a single or \textit{scalar} conservation law (i.e., $m=1$ above).

Following then the development in \S3.4.1, let us suppose

$$(12) \qquad \begin{cases} \mathbf{v}: \mathbb{R}\times[0,\infty) \to \mathbb{R}^m \text{ is smooth,} \\ \text{with compact support, } \mathbf{v}=(v^1,\ldots,v^m). \end{cases}$$

We temporarily assume $\mathbf{u}$ is a smooth solution of our problem (5), take the dot product of the PDE $\mathbf{u}_t + \mathbf{F}(\mathbf{u})_x = 0$ with the test function $\mathbf{v}$, and integrate by parts, to obtain the equality

$$(13) \qquad \int_0^\infty \int_{-\infty}^\infty \mathbf{u}\cdot\mathbf{v}_t + \mathbf{F}(\mathbf{u})\cdot\mathbf{v}_x \, dx\,dt + \int_{-\infty}^\infty \mathbf{g}\cdot\mathbf{v}\,dx\Big|_{t=0} = 0.$$

This identity, which we derived supposing $\mathbf{u}$ to be a smooth solution, makes sense if $\mathbf{u}$ is merely bounded.

\begin{definition}[Integral solution of a system of conservation laws]
\label{def:integral-solution-system}
\dcref{ch11:11.1-def-1}
\group{Conservation Laws}
\uses{def:integral-solution-conservation-law}
We say that $\mathbf{u}\in L^\infty(\mathbb{R}^n\times(0,\infty);\mathbb{R}^m)$ is an \textit{integral solution} of the initial-value problem (5) provided the equality (13) holds for all test functions $\mathbf{v}$ satisfying (12).
\end{definition}

Continuing now to parallel the development in \S3.4.1 for a single conservation law, let us now consider the situation that we have an integral solution $\mathbf{u}$ of (5) which is smooth on either side of a curve $C$, along which $\mathbf{u}$ has simple jump discontinuities. More precisely, let us assume that $V \subset \mathbb{R}\times(0,\infty)$ is some region cut by a smooth curve $C$ into a left hand part $V_l$ and a right hand part $V_r$.

\begin{figure}[h]\centering
% [FIGURE: A region V in the (x,t)-plane divided by a smooth curve C into a left-hand shaded region V_l and a right-hand shaded region V_r; caption "Rankine-Hugoniot condition"]
\end{figure}

\begin{proposition}[Rankine--Hugoniot jump condition]
\label{prop:rankine-hugoniot-condition-system}
\dcref{ch11:11.1-prop-1}
\group{Conservation Laws}
\uses{def:integral-solution-system}
Suppose the curve $C$ is represented parametrically as $\{(x,t)\mid x = s(t)\}$ for some smooth function $s(\cdot):[0,\infty)\to\mathbb{R}$. Then, along $C$,
$$(20) \qquad [\![\mathbf{F}(\mathbf{u})]\!] = \sigma[\![\mathbf{u}]\!],$$
where
$$\begin{cases}
[\![\mathbf{u}]\!] = \mathbf{u}_l - \mathbf{u}_r = \text{jump in } \mathbf{u} \text{ across the curve } C \\
[\![\mathbf{F}(\mathbf{u})]\!] = \mathbf{F}(\mathbf{u}_l) - \mathbf{F}(\mathbf{u}_r) = \text{jump in } \mathbf{F}(\mathbf{u}) \\
\sigma = \dot{s} = \text{speed of the curve } C,
\end{cases}$$
the subscript ``$l$'' denoting the limit from the left and ``$r$'' the limit from the right. Note carefully this is a vector equality.
\end{proposition}

\begin{proof}
Assuming that $\mathbf{u}$ is smooth in $V_l$, we select the test function $\mathbf{v}$ with compact support in $V_l$ and deduce from (13) that
$$(14) \qquad \mathbf{u}_t + \mathbf{F}(\mathbf{u})_x = 0 \quad \text{in } V_l.$$
Similarly, we have
$$(15) \qquad \mathbf{u}_t + \mathbf{F}(\mathbf{u})_x = 0 \quad \text{in } V_r,$$
provided $\mathbf{u}$ is smooth in $V_r$. Now choose a test function $\mathbf{v}$ with compact support in $V$, but which does not necessarily vanish along the curve $C$. Then utilizing the identity (13) we find
$$(16) \qquad 0 = \int_0^\infty\int_{-\infty}^\infty \mathbf{u}\cdot\mathbf{v}_t + \mathbf{F}(\mathbf{u})\cdot\mathbf{v}_x\,dx\,dt = \iint_{V_l}\mathbf{u}\cdot\mathbf{v}_t+\mathbf{F}(\mathbf{u})\cdot\mathbf{v}_x\,dx\,dt + \iint_{V_r}\mathbf{u}\cdot\mathbf{v}_t+\mathbf{F}(\mathbf{u})\cdot\mathbf{v}_x\,dx\,dt.$$
As $\mathbf{v}$ has compact support within $V$, we deduce
$$(17) \qquad \iint_{V_l}\mathbf{u}\cdot\mathbf{v}_t+\mathbf{F}(\mathbf{u})\cdot\mathbf{v}_x\,dx\,dt = -\iint_{V_l}[\mathbf{u}_t+\mathbf{F}(\mathbf{u})_x]\cdot\mathbf{v}\,dx\,dt + \int_C(\mathbf{u}_l\nu^2+\mathbf{F}(\mathbf{u}_l)\nu^1)\cdot\mathbf{v}\,dl = \int_C(\mathbf{u}_l\nu^2+\mathbf{F}(\mathbf{u}_l)\nu^1)\cdot\mathbf{v}\,dl,$$
owing to (14). Here $\nu=(\nu^1,\nu^2)$ is the unit normal to the curve $C$ pointing from $V_l$ into $V_r$, and the subscript ``$l$'' denotes the limit from the left. Likewise (15) implies
$$\iint_{V_r}\mathbf{u}\cdot\mathbf{v}_t+\mathbf{F}(\mathbf{u})\cdot\mathbf{v}_x\,dx\,dt = -\int_C(\mathbf{u}_r\nu^2+\mathbf{F}(\mathbf{u}_r)\nu^1)\cdot\mathbf{v}\,dl,$$
``$r$'' denoting the limit from the right. Adding this identity to (17) and remembering (16), we deduce
$$\int_C[(\mathbf{F}(\mathbf{u}_l)-\mathbf{F}(\mathbf{u}_r))\nu^1 + (\mathbf{u}_l-\mathbf{u}_r)\nu^2]\cdot\mathbf{v}\,dl = 0.$$
This identity obtains for all smooth functions $\mathbf{v}$ as above; whence
$$(18) \qquad (\mathbf{F}(\mathbf{u}_l)-\mathbf{F}(\mathbf{u}_r))\nu^1 + (\mathbf{u}_l-\mathbf{u}_r)\nu^2 = 0 \quad \text{along } C.$$
Suppose now the curve $C$ is represented parametrically as $\{(x,t)\mid x=s(t)\}$ for some smooth function $s(\cdot):[0,\infty)\to\mathbb{R}$. Then $\nu=(\nu^1,\nu^2)=(1+\dot s^2)^{-1/2}(1,-\dot s)$. Consequently (18) reads
$$(19) \qquad \mathbf{F}(\mathbf{u}_l)-\mathbf{F}(\mathbf{u}_r) = \dot s(\mathbf{u}_l-\mathbf{u}_r)$$
in $V$ along the curve $C$. Rewriting (19) using the jump/speed notation above gives (20), the \textit{Rankine--Hugoniot jump condition}.
\end{proof}

\begin{figure}[h]\centering
% [FIGURE: A small empty square box]
\end{figure}

\subsection{Traveling waves, hyperbolic systems}\label{sec:traveling-waves-hyperbolic-systems}

We have seen in \S3.4 that the notion of an integral solution for conservation laws is not adequate: such solutions need not be unique. We are therefore intent upon discovering some additional requirements for a good definition of a generalized solution. This will presumably entail as in \S3.4 an entropy criterion based upon an analysis of shock waves. This expectation, now as carried over to systems, is largely correct, but first of all we must study more carefully the nonlinearity $F$ in the hopes of discovering mathematically appropriate and physically correct structural conditions to impose.

Let us start by first considering the wider class of \textit{semilinear systems} having the nondivergence form:

$$(21) \qquad \mathbf{u}_t + \mathbf{B}(\mathbf{u})\mathbf{u}_x = 0 \quad \text{in } \mathbb{R}\times(0,\infty),$$

where $\mathbf{B}:\mathbb{R}^m\to M^{m\times m}$. This system is for smooth functions equivalent to the conservation law in (5), provided

$$\mathbf{B} = DF = \begin{pmatrix} F_{z_1}^1 & \cdots & F_{z_m}^1 \\ \vdots & \ddots & \vdots \\ F_{z_1}^m & \cdots & F_{z_m}^m \end{pmatrix}_{m\times m}.$$

We consider now the possibility of finding particular solutions which have the form of a traveling wave:

$$(22) \qquad \mathbf{u}(x,t) = \mathbf{v}(x-\sigma t) \quad (x\in\mathbb{R}, t>0),$$

where the profile $\mathbf{v}:\mathbb{R}\to\mathbb{R}^m$ and the velocity $\sigma\in\mathbb{R}$ are to be found. We substitute expression (22) into the PDE (21), and thereby obtain the equality

$$(23) \qquad -\sigma\mathbf{v}'(x-\sigma t) + \mathbf{B}(\mathbf{v}(x-\sigma t))\mathbf{v}'(x-\sigma t) = 0.$$

Observe (23) says that $\sigma$ is an eigenvalue of the matrix $\mathbf{B}(\mathbf{v})$ corresponding to the eigenvector $\mathbf{v}'$.

This conclusion suggests (exactly as for the linear theory in \S7.3) that if we wish to find traveling waves or, more generally, wavelike solutions of our system of PDE, we should make some sort of hyperbolicity hypothesis concerning the eigenvalues of $\mathbf{B}$.

\begin{definition}[Strict hyperbolicity]
\label{def:strictly-hyperbolic-system}
\dcref{ch11:11.1-def-2}
\group{Conservation Laws}
If for each $z\in\mathbb{R}^m$ the eigenvalues of $\mathbf{B}(z)$ are real and distinct, we call the system (21) \textit{strictly hyperbolic}.
\end{definition}

We henceforth assume the system of partial differential equations (21) (and the special case $\mathbf{B}=DF$ of conservation laws) to be strictly hyperbolic.

\begin{definition}[Eigenvalues and eigenvectors of a hyperbolic system]
\label{def:eigenvalue-eigenvector-notation-hyperbolic}
\dcref{ch11:11.1-def-3}
\group{Conservation Laws}
\uses{def:strictly-hyperbolic-system}
(i) We will write
$$(24) \qquad \lambda_1(z) < \lambda_2(z) < \cdots < \lambda_m(z) \quad (z\in\mathbb{R}^m)$$
to denote the real and distinct eigenvalues of $\mathbf{B}(z)$, in increasing order.

(ii) Then for each $k=1,\ldots,m$, we let
$$\mathbf{r}_k(z)$$
denote a corresponding nonzero eigenvector, so that
$$(25) \qquad \mathbf{B}(z)\mathbf{r}_k(z) = \lambda_k(z)\mathbf{r}_k(z) \quad (k=1,\ldots,m,\ z\in\mathbb{R}^m).$$
Since we are always assuming the strict hyperbolicity condition, the vectors $\{\mathbf{r}_k(z)\}_{k=1}^m$ span $\mathbb{R}^m$ for each $z\in\mathbb{R}^m$.

(iii) Next, since a matrix and its transpose have the same spectrum, we can introduce for each $k=1,\ldots,m$ a nonzero eigenvector
$$\mathbf{l}_k(z)$$
for the matrix $\mathbf{B}(z)^T$, corresponding to the eigenvalue $\lambda_k(z)$. Thus
$$(26) \qquad \mathbf{B}(z)^T\mathbf{l}_k(z) = \lambda_k(z)\mathbf{l}_k(z) \quad (k=1,\ldots,m,\ z\in\mathbb{R}^m).$$
This equality is usually written
$$(27) \qquad \mathbf{l}_k(z)\mathbf{B}(z) = \lambda_k(z)\mathbf{l}_k(z) \quad (k=1,\ldots,m,\ z\in\mathbb{R}^m).$$
Thus $\{\mathbf{l}_k(z)\}_{k=1}^m$ can be regarded as \textit{left eigenvectors} of $\mathbf{B}(z)$, and $\{\mathbf{r}_k(z)\}_{k=1}^m$ are \textit{right eigenvectors}. $\square$
\end{definition}

\begin{remark}[Orthogonality of left and right eigenvectors]
\label{rem:eigenvector-orthogonality}
\dcref{ch11:11.1-rem-1}
\group{Conservation Laws}
\uses{def:eigenvalue-eigenvector-notation-hyperbolic}
Additionally, we observe
$$(28) \qquad \mathbf{l}_l(z)\cdot\mathbf{r}_k(z) = 0 \quad \text{if } k\neq l \quad (z\in\mathbb{R}^m).$$
To confirm this, we compute using (25) and (26) that
$$\lambda_k(z)(\mathbf{l}(z)\cdot\mathbf{r}_k(z)) = \mathbf{l}(z)\cdot(\mathbf{B}(z)\mathbf{r}_k(z)) = (\mathbf{B}(z)^T\mathbf{l}(z))\cdot\mathbf{r}_k(z) = \lambda_l(z)(\mathbf{l}(z)\cdot\mathbf{r}_k(z));$$
whence (28) follows since $\lambda_k(z)\neq\lambda_l(z)$ if $k\neq l$. $\square$
\end{remark}

Let us first show that the notion of strict hyperbolicity is independent of coordinates.

\begin{theorem}[Invariance of hyperbolicity under change of coordinates]
\label{thm:invariance-hyperbolicity-coordinate-change}
\dcref{ch11:11.1-thm-1}
\group{Conservation Laws}
\uses{def:strictly-hyperbolic-system,def:eigenvalue-eigenvector-notation-hyperbolic}
Let $u$ be a smooth solution of the strictly hyperbolic system (21). Assume also $\Phi:\mathbb{R}^m\to\mathbb{R}^m$ is a smooth diffeomorphism, with inverse $\Psi$. Then
$$(29) \qquad \tilde{u} := \Phi(u)$$
solves the strictly hyperbolic system
$$(30) \qquad \tilde u_t + \tilde B(\tilde u)\tilde u_x = 0 \quad \text{in } \mathbb{R}\times(0,\infty),$$
for
$$(31) \qquad \tilde B(z) := D\Phi(\Psi(z))B(\Psi(z))D\Psi(z) \quad (z\in\mathbb{R}^m).$$
\end{theorem}

\begin{proof}
1. We compute $\tilde u_t = D\Phi(u)u_t$, $\tilde u_x = D\Phi(u)u_x$, and so equation (30) is valid for $\tilde B(z) = D\Phi(z)B(z)D\Phi^{-1}(z)$, where $z=\Phi(z)$. Substituting $z=\Psi(\tilde z)$, we obtain (31).

2. We must prove the system (30) is strictly hyperbolic. If $\lambda_k(z)$ is an eigenvalue of $B(z)$, with corresponding right eigenvector $r_k(z)$, we have
$$B(z)r_k(z) = \lambda_k(z)r_k(z).$$
Setting
$$(32) \qquad \tilde r_k(\tilde z) := D\Phi(\Psi(\tilde z))r_k(\Psi(\tilde z)),$$
$$(33) \qquad \tilde\lambda_k(\tilde z) := \lambda_k(\Psi(\tilde z)),$$
we compute
$$(34) \qquad \tilde B(\tilde z)\tilde r_k(\tilde z) = \tilde\lambda_k(\tilde z)\tilde r_k(\tilde z).$$
Similarly if $l_k(z)$ is a left eigenvector, we write
$$(35) \qquad \tilde l_k(\tilde z) := l_k(\Psi(\tilde z))D\Psi(\tilde z),$$
and calculate
$$(36) \qquad \tilde l_k(\tilde z)\tilde B(z) = \tilde\lambda_k(\tilde z)\tilde l_k(\tilde z).$$
In view of (32)--(36), we conclude that the system (30) is strictly hyperbolic. $\square$
\end{proof}

Next we study how $\lambda_k(z)$, $r_k(z)$ and $l_k(z)$ change as $z$ varies:

\begin{theorem}[Dependence of eigenvalues and eigenvectors on parameters]
\label{thm:smooth-dependence-eigenvalues-eigenvectors}
\dcref{ch11:11.1-thm-2}
\group{Conservation Laws}
\uses{def:strictly-hyperbolic-system,def:eigenvalue-eigenvector-notation-hyperbolic}
Assume the matrix function $\mathbf{B}$ is smooth, strictly hyperbolic. Then

(i) the eigenvalues $\lambda_k(z)$ depend smoothly on $z\in\mathbb{R}^m$ $(k=1,\ldots,m)$.

(ii) Furthermore, we can select the right eigenvectors $\mathbf{r}_k(z)$ and left eigenvectors $\mathbf{l}_k(z)$ to depend smoothly on $z\in\mathbb{R}^m$ and satisfy the normalization
$$|\mathbf{r}_k(z)|,\ |\mathbf{l}_k(z)| = 1 \quad (k=1,\ldots,m).$$
\end{theorem}

\begin{proof}
1. Since $\mathbf{B}(z)$ is strictly hyperbolic, for each $z_0\in\mathbb{R}^m$ we have
$$(37) \qquad \lambda_1(z_0) < \lambda_2(z_0) < \cdots < \lambda_m(z_0).$$
Fix $k\in\{1,\ldots,m\}$ and any point $z_0\in\mathbb{R}^m$, and let $\mathbf{r}_k(z_0)$ satisfy
$$\begin{cases} B(z_0)\mathbf{r}_k(z_0) = \lambda_k(z_0)\mathbf{r}_k(z_0) \\ |\mathbf{r}_k(z_0)| = 1. \end{cases}$$
Upon rotating coordinates if necessary, we may assume
$$(38) \qquad \mathbf{r}_k(z_0) = e_m = (0,\ldots,1).$$
We first show that near $z_0$, there exist smooth functions $\lambda_k(z)$, $\mathbf{r}_k(z)$ such that
$$\begin{cases} B(z)\mathbf{r}_k(z) = \lambda_k(z)\mathbf{r}_k(z) \\ |\mathbf{r}_k(z)| = 1. \end{cases}$$

2. We will apply the Implicit Function Theorem to the smooth function $\Phi:\mathbb{R}^m\times\mathbb{R}\times\mathbb{R}^m\to\mathbb{R}^{m+1}$ defined by
$$\Phi(r,\lambda,z) = (\mathbf{B}(z)r - \lambda r, |r|^2) \quad (r,z\in\mathbb{R}^m,\ \lambda\in\mathbb{R}).$$
Now
$$\frac{\partial\Phi(r,\lambda,z)}{\partial(r,\lambda)} = \begin{pmatrix} \mathbf{B}(z)-\lambda I & \vdots & -r_1 \\ & \vdots & \vdots \\ & \vdots & -r_m \\ 2r_1 \ldots 2r_m & & 0 \end{pmatrix}_{(m+1)\times(m+1)},$$
and so, according to (38), it suffices to check that
$$(39) \qquad \det\begin{pmatrix} \mathbf{B}(z_0)-\lambda_k(z_0)I & \vdots & 0 \\ & \vdots & \vdots \\ & \vdots & -1 \\ 0\ldots\ldots 2 & & 0 \end{pmatrix} \neq 0.$$

3. Note that for $\epsilon>0$ sufficiently small, the matrix

$$(40) \qquad \mathbf{B}_\epsilon = B(z_0) - (\lambda_k(z_0)+\epsilon)I$$

is invertible. In light of (38),
$$\mathbf{B}_\epsilon e_m = -\epsilon e_m.$$
Therefore
$$\begin{pmatrix} \mathbf{B}_\epsilon & \vdots & 0 \\ -1 & \vdots & \\ 0\ldots 2 & 0 \end{pmatrix} \begin{pmatrix} I & \vdots & 0 \\ & (-\epsilon)^{-1} & \\ 0\ldots 0 & 1 \end{pmatrix} = \begin{pmatrix} \mathbf{B}_\epsilon & \vdots & 0 \\ & & 0 \\ 0\ldots 2 & 2(-\epsilon)^{-1} \end{pmatrix}.$$
Consequently, since the determinant of the second matrix before the equals sign is one, we have
$$\det\begin{pmatrix} \mathbf{B}_\epsilon & \vdots & 0 \\ -1 & \vdots & \\ 0\ldots 2 & 0 \end{pmatrix} = 2(\det\mathbf{B}_\epsilon)(-\epsilon)^{-1} = 2\prod_{j\neq k}(\lambda_j(z_0)-(\lambda_k(z_0)+\epsilon))(-\epsilon)(-\epsilon)^{-1} \to 2\prod_{j\neq k}(\lambda_j(z_0)-\lambda_k(z_0)) \quad \text{as } \epsilon\to 0.$$
As $B(z_0)$ is strictly hyperbolic, the last expression is nonzero. Condition (39) is verified. We may thus invoke the Implicit Function Theorem to find near $z_0$ smooth functions $\lambda_k(z)$ and $\mathbf{r}_k(z)$, satisfying the conclusion of the theorem.

4. It remains to show that we can define $\lambda_k(z)$ and $\mathbf{r}_k(z)$ for all $z\in\mathbb{R}^m$, and not just near any particular point $z_0$. To do so, let us write
$$R := \sup\{r>0 \mid \lambda_k(z),\mathbf{r}_k(z) \text{ as above exist and are smooth on } B(0,r)\}.$$
If $R=\infty$, we are done. Otherwise, we cover $\partial B(0,R)$ with finitely many open balls into which we can smoothly extend $\lambda_k(\cdot)$ and $\mathbf{r}_k(\cdot)$, using steps 1--3 above. This yields a contradiction to the definition of $R$.

A similar proof works for the left eigenvectors. $\square$
\end{proof}

\begin{remark}[Global orientation of eigenspaces]
\label{rem:global-orientation-eigenspaces}
\dcref{ch11:11.1-rem-2}
\group{Conservation Laws}
\uses{thm:smooth-dependence-eigenvalues-eigenvectors}
Observe that we are not only globally and smoothly defining the eigenvalues and eigenspaces of $\mathbf{B}$, but are also globally providing the eigenspaces with an orientation.

This proof depends fundamentally upon the one-dimensionality of the eigenspaces. See Problem 2 for an example of what could go wrong in the hyperbolic, but not strictly hyperbolic, setting. $\square$
\end{remark}

\begin{example}[Strict hyperbolicity of the $p$-system]
\label{ex:p-system-strict-hyperbolicity}
\dcref{ch11:11.1-ex-4}
\group{Conservation Laws}
\uses{ex:p-system-definition,def:strictly-hyperbolic-system}
For the $p$-system (6), we have
$$\mathbf{u}_t + \mathbf{B}(\mathbf{u})\mathbf{u}_x = 0,$$
for
$$\mathbf{B}(z) = DF(z) = \begin{pmatrix} 0 & -1 \\ -p'(z_1) & 0 \end{pmatrix}.$$
The eigenvalues are $\lambda_1=-\sigma$, $\lambda_2=\sigma$, for $\sigma:=p'(z_1)^{1/2}$. These are real and distinct provided we hereafter suppose the strict hyperbolicity condition
$$(41) \qquad p'>0.$$
For the nonlinear wave equation (8) this is the physical assumption that the \textit{stress} $p(u_x)$ is a strictly increasing function of the \textit{strain} $u_x$. $\square$
\end{example}

\begin{example}[Strict hyperbolicity of Euler's equations]
\label{ex:euler-equations-strict-hyperbolicity}
\dcref{ch11:11.1-ex-5}
\group{Conservation Laws}
\uses{ex:euler-equations-1d-system,thm:invariance-hyperbolicity-coordinate-change,def:strictly-hyperbolic-system}
Euler's equations (9) comprise a strictly hyperbolic system provided we assume $p>0$ and
$$(42) \qquad \frac{\partial p}{\partial\rho}>0, \quad \frac{\partial p}{\partial e}>0,$$
where $p=p(\rho,e)$ is the constitutive relation between the mass density, the internal energy density and the pressure. This assertion is however difficult to verify directly, as the flux function $\mathbf{F}$ defined by (11) is complicated.

Let us rather change variables and regard the density $\rho$, velocity $v$ and internal energy $e$ as the unknowns. We can then rewrite Euler's equations (9) in terms of these quantities, and, so doing, obtain after some calculations the system
$$(43) \qquad \begin{cases} \rho_t + v\rho_x + \rho v_x = 0 \\ v_t + vv_x + \frac{1}{\rho}p_x = 0 \\ e_t + ve_x + \frac{v}{\rho}v_x = 0, \end{cases}$$
provided $\rho>0$. These equations are not in conservation form. Setting now $\mathbf{u}=(u^1,u^2,u^3)=(\rho,v,e)$, we rewrite (43) as
$$(44) \qquad \mathbf{u}_t + \mathbf{B}(\mathbf{u})\mathbf{u}_x = 0 \quad \text{in } \mathbb{R}\times(0,\infty),$$
for
$$(45) \qquad \mathbf{B}(z) = z_2 I + \tilde{\mathbf{B}}(z),$$
where
$$\tilde{\mathbf{B}}(z) := \begin{pmatrix} 0 & z_1 & 0 \\ \frac{1}{z_1}\frac{\partial p}{\partial\rho}(z_1,z_3) & 0 & \frac{1}{z_1}\frac{\partial p}{\partial\varepsilon}(z_1,z_3) \\ 0 & \frac{1}{z_1}p(z_1,z_3) & 0 \end{pmatrix}.$$
The characteristic polynomial of $\mathbf{B}$ is $-\lambda(\lambda^2-\sigma^2)$, for $\sigma^2 = \frac{p}{\rho^2}\frac{\partial p}{\partial\varepsilon} + \frac{\partial p}{\partial\rho}$. Recalling (45) and reverting to physical notation, we see that the eigenvalues of $\mathbf{B}$ are
$$(46) \qquad \lambda_1 = v-\sigma, \quad \lambda_2 = v, \quad \lambda_3 = v+\sigma,$$
where
$$\sigma := \left(\frac{p}{\rho^2}\frac{\partial p}{\partial\varepsilon} + \frac{\partial p}{\partial\rho}\right)^{1/2} > 0$$
is the local \textit{sound speed}. We therefore see that the system (44) is strictly hyperbolic, provided assumption (42) is valid. Remembering now \cref{thm:invariance-hyperbolicity-coordinate-change}, we deduce that Euler's equations (9) are also strictly hyperbolic, with eigenvalues given by (46). $\square$
\end{example}

\section{Riemann's Problem}\label{sec:riemanns-problem}

In this section we investigate in detail the system of conservation laws

$$(1) \qquad \mathbf{u}_t + \mathbf{F}(\mathbf{u})_x = 0 \quad \text{in } \mathbb{R}\times(0,\infty),$$

with the piecewise-constant initial data

$$(2) \qquad \mathbf{g} = \begin{cases} \mathbf{u}_l & \text{if } x<0 \\ \mathbf{u}_r & \text{if } x>0. \end{cases}$$

This is \textit{Riemann's problem}. We call the given vectors $\mathbf{u}_l$ and $\mathbf{u}_r$ the left and right \textit{initial states}.

\subsection{Simple waves}\label{sec:simple-waves}

We commence our study of (1), (2) in very much the same spirit as in \S11.1.2, in that we look for solutions of (1) having a special form. Before we searched for traveling waves, that is, solutions of the type $\mathbf{u}(x,t) = \mathbf{v}(x-\sigma t)$. We now seek \textit{simple waves}. These are solutions of (1) having the structure

$$(3) \qquad \mathbf{u}(x,t) = \mathbf{v}(w(x,t)) \quad (x\in\mathbb{R}, t>0),$$

where $\mathbf{v}:\mathbb{R}\to\mathbb{R}^m$, $\mathbf{v}=(v^1,\ldots,v^m)$, and $w:\mathbb{R}\times[0,\infty)\to\mathbb{R}$ are to be found. To discover the requisite properties of $\mathbf{v}$ and $w$, let us substitute (3) into (1) and obtain the equality

$$(4) \qquad \dot{\mathbf{v}}(w) + DF(\mathbf{v}(w))\dot{\mathbf{v}}(w)w_x = 0.$$

Now in view of equation (25) from \S11.1.2, with $B=DF$, we see (4) will be valid if for some $k\in\{1,\ldots,m\}$ $w$ solves the PDE

$$(5) \qquad w_t + \lambda_k(\mathbf{v}(w))w_x = 0$$

and $\mathbf{v}$ solves the ODE

$$(6) \qquad \dot{\mathbf{v}}(s) = \mathbf{r}_k(\mathbf{v}(s)) \quad \left(\cdot = \frac{d}{ds}\right).$$

If (5) and (6) hold, we call the function $u$ defined by (3) a $k$-\textit{simple wave}. The point of all this is that we can regard (6) as an ODE for the vector function $\mathbf{v}$, and then---once $\mathbf{v}$ has been found by solving (6)---can interpret equation (5) as a \textit{scalar} conservation law for $w$.

Let us next identify circumstances under which we can employ the construction (3)--(6) to build a continuous solution $\mathbf{u}$ of (1). We must examine first the ODE (6).

\begin{definition}[Rarefaction curve]
\label{def:rarefaction-curve}
\dcref{ch11:11.2-def-1}
\group{Conservation Laws}
\uses{def:eigenvalue-eigenvector-notation-hyperbolic}
Given a fixed state $z_0\in\mathbb{R}^m$, we define the $k^{\text{th}}$-\textit{rarefaction curve}
$$R_k(z_0)$$
to be the path in $\mathbb{R}^m$ of the solution of the ODE (6) which passes through $z_0$.
\end{definition}

\begin{figure}[h]\centering
% [FIGURE: A curve labeled $\mathbf{R}_k(\mathbf{z}_0)$ passing through a point labeled $z_0$, with arrows on both ends indicating the curve extends in both directions. The curve has a characteristic U or valley-like shape. Below the figure is a caption "Rarefaction curve"]
\end{figure}

Given then the solution $\mathbf{v}$ of (6), we turn to the PDE (5), which we rewrite as the scalar conservation law

$$(7) \qquad w_t + F_k(w)_x = 0$$

for

$$(8) \qquad F_k(s) := \int_0^s \lambda_k(v(t))\,dt \quad (s\in\mathbb{R}).$$

The PDE (7) will fall under the general theory developed in \S3.4 provided $F_k$ is strictly convex (or else strictly concave). Let us therefore compute

$$(9) \qquad F_k'(s) = \lambda_k(v(s)),$$

$$(10) \qquad F_k''(s) = D\lambda_k(v(s))\cdot\dot v(s) = D\lambda_k(v(s))\cdot r_k(v(s)).$$

Owing to (10), the function $F_k$ will be convex if
$$D\lambda_k(z)\cdot r_k(z) > 0 \quad (z\in\mathbb{R}^m),$$
and concave if
$$D\lambda_k(z)\cdot r_k(z) < 0 \quad (z\in\mathbb{R}^m).$$
The function $F_k$ is linear provided
$$D\lambda_k(z)\cdot r_k(z) \equiv 0 \quad (z\in\mathbb{R}^m).$$

These possibilities motivate the following

\begin{definition}[Genuine nonlinearity and linear degeneracy]
\label{def:genuinely-nonlinear-linearly-degenerate}
\dcref{ch11:11.2-def-2}
\group{Conservation Laws}
\uses{def:eigenvalue-eigenvector-notation-hyperbolic}
(i) The pair $(\lambda_k(z), r_k(z))$ is called \textit{genuinely nonlinear} provided
$$(11) \qquad D\lambda_k(z)\cdot r_k(z) \neq 0 \quad \text{for all } z\in\mathbb{R}^m.$$

(ii) We say $(\lambda_k(z), r_k(z))$ is \textit{linearly degenerate} if
$$(12) \qquad D\lambda_k(z)\cdot r_k(z) = 0 \quad \text{for all } z\in\mathbb{R}^m.$$
\end{definition}

\textbf{Notation.} If the pair $(\lambda_k, r_k)$ is genuinely nonlinear, write
$$R_k^+(z_0) := \{z\in R_k(z_0) \mid \lambda_k(z) > \lambda_k(z_0)\}$$
and
$$R_k^-(z_0) := \{z\in R_k(z_0) \mid \lambda_k(z) < \lambda_k(z_0)\}.$$
Then
$$R_k(z_0) = R_k^+(z_0) \cup \{z_0\} \cup R_k^-(z_0).$$

\begin{figure}[h]\centering
% [FIGURE: A diagram showing the rarefaction curve split into two parts at a point $z_0$. The left part is labeled $R_k^-(z_0)$ and curves upward from the left toward $z_0$. The right part is labeled $R_k^+(z_0)$ and curves upward from $z_0$ toward the right. Both parts meet at the point $z_0$ on the horizontal axis.]
\end{figure}

Two parts of the rarefaction curve

\subsection{Rarefaction waves}\label{sec:rarefaction-waves}

We turn our attention again to Riemann's problem (1), (2).

\begin{theorem}[Existence of $k$-rarefaction waves]
\label{thm:existence-k-rarefaction-waves}
\dcref{ch11:11.2-thm-1}
\group{Conservation Laws}
\uses{def:genuinely-nonlinear-linearly-degenerate,def:rarefaction-curve,thm:riemann-problem-solution}
Suppose that for some $k\in\{1,\ldots,m\}$,

(i) the pair $(\lambda_k, r_k)$ is genuinely nonlinear, and

(ii) $u_r \in R_k^+(u_l)$.

Then there exists a continuous integral solution $\mathbf{u}$ of Riemann's problem (1), (2), which is a $k$-simple wave constant along lines through the origin.
\end{theorem}

\begin{remark}[Centered $k$-rarefaction wave]
\label{rem:k-rarefaction-wave-terminology}
\dcref{ch11:11.2-rem-1}
\group{Conservation Laws}
\uses{thm:existence-k-rarefaction-waves}
We call $\mathbf{u}$ a (centered) $k$-\textit{rarefaction wave}.
\end{remark}

\begin{figure}[h]\centering
% [FIGURE: A diagram showing a k-rarefaction wave with multiple line segments emanating from a central point on a horizontal axis, fanning outward and upward to the left and right, illustrating how characteristics spread out from the origin]
\end{figure}

\begin{proof}
1. We first choose $w_l$ and $w_r \in \mathbb{R}$ so that $u_l = \mathbf{v}(w_l)$, $u_r = \mathbf{v}(w_r)$. Suppose for the moment
$$(13) \qquad w_l < w_r.$$

2. Consider then the scalar Riemann problem consisting of the PDE (7) together with the initial condition
$$(14) \qquad g = \begin{cases} w_l & \text{if } x<0 \\ w_r & \text{if } x>0. \end{cases}$$
Now in view of hypothesis (ii) we have $\lambda_k(u_r) > \lambda_k(u_l)$; that is, according to (9), $F_k'(w_r) > F_k'(w_l)$. But then it follows from (i) that the function $F_k$ defined by (8) is strictly convex. Accordingly we can apply Theorem 4 in \S3.4.4 to the scalar Riemann problem (7), (14), whose unique weak solution is a continuous rarefaction wave connecting the states $w_l$ and $w_r$. More specifically,
$$w(x,t) = \begin{cases} w_l & \text{if } \frac{x}{t} < F_k'(w_l) \\ G_k\left(\frac{x}{t}\right) & \text{if } F_k'(w_l) < \frac{x}{t} < F_k'(w_r) \quad (x\in\mathbb{R}, t>0) \\ w_r & \text{if } \frac{x}{t} > F_k'(w_r), \end{cases}$$
where $G_k = (F_k')^{-1}$. Thus $\mathbf{u}(x,t) = \mathbf{v}(w(x,t))$, where $\mathbf{v}$ solves the ODE (6) and passes through $u_l$, is a continuous integral solution of (1), (2).

3. The case $w_l > w_r$ is treated similarly, since $F_k$ is then concave. $\square$
\end{proof}

\subsection{Shock waves, contact discontinuities}\label{sec:shock-waves-contact-discontinuities}

We consider next the possibility that the states $u_l$ and $u_r$ may be joined not by a rarefaction wave as above, but rather by a shock.

\begin{figure}[h]\centering
% [FIGURE: A diagram showing a k-shock wave: a solid vertical line rising from the origin with characteristics represented as diagonal lines converging into the shock from both sides; caption "k-shock wave"]
\end{figure}

\subsubsection*{a. The shock set}

Recalling the Rankine--Hugoniot condition from \S11.1.1, we see that necessarily we must have the equality $\mathbf{F}(u_l) - \mathbf{F}(u_r) = \sigma(u_l - u_r)$, where $\sigma\in\mathbb{R}$, for such a shock wave to exist. This observation motivates the following

\begin{definition}[Shock set]
\label{def:shock-set}
\dcref{ch11:11.2-def-3}
\group{Conservation Laws}
\uses{prop:rankine-hugoniot-condition-system}
Given a fixed state $z_0\in\mathbb{R}^m$, we define the \textit{shock set}
$$S(z_0) := \{z\in\mathbb{R}^m \mid \mathbf{F}(z) - \mathbf{F}(z_0) = \sigma(z-z_0) \text{ for a constant } \sigma = \sigma(z,z_0)\}.$$
\end{definition}

\begin{theorem}[Structure of the shock set]
\label{thm:structure-of-shock-set}
\dcref{ch11:11.2-thm-2}
\group{Conservation Laws}
\uses{def:shock-set,thm:smooth-dependence-eigenvalues-eigenvectors,def:eigenvalue-eigenvector-notation-hyperbolic}
Fix $z_0\in\mathbb{R}^m$. In some neighborhood of $z_0$, $S(z_0)$ consists of the union of $m$ smooth curves $S_k(z_0)$ $(k=1,\ldots,m)$, with the following properties:

(i) The curve $S_k(z_0)$ passes through $z_0$, with tangent $\mathbf{r}_k(z_0)$.

(ii) $\displaystyle\lim_{\substack{z\to z_0 \\ z\in S_k(z_0)}} \sigma(z,z_0) = \lambda_k(z_0)$.

(iii) $\sigma(z,z_0) = \dfrac{\lambda_k(z)+\lambda_k(z_0)}{2} + O(|z-z_0|^2)$, as $z\to z_0$ with $z\in S_k(z_0)$.
\end{theorem}

\begin{proof}
1. Define
$$\mathbf{B}(z) := \int_0^1 DF(z_0+t(z-z_0))\,dt \quad (z\in\mathbb{R}^m).$$
Then
$$(15) \qquad \mathbf{B}(z)(z-z_0) = \mathbf{F}(z) - \mathbf{F}(z_0).$$
In particular $z\in S(z_0)$ if and only if
$$(16) \qquad (\mathbf{B}(z)-\sigma I)(z-z_0) = 0$$
for some scalar $\sigma = \sigma(z,z_0)$.

\begin{figure}[h]\centering
% [FIGURE: A curve diagram showing contact between $R_k$ and $S_k$. The diagram depicts two curved trajectories meeting at a point labeled $z_0$. The left curve is labeled $R_k(z_0)$ and the right curve is labeled $S_k(z_0)$. Multiple arrows point outward from the curves indicating flow direction. The curves make smooth contact at $z_0$.]
\end{figure}

Contact between $R_k$ and $S_k$

2. We study equation (16) by first of all noting
$$(17) \qquad \mathbf{B}(z_0) = D\mathbf{F}(z_0).$$
Now in view of the strict hyperbolicity, the characteristic polynomial $\lambda \mapsto \det(\lambda I - \mathbf{B}(z_0))$ has $m$ distinct, real roots, and hence the polynomial $\lambda \mapsto \det(\lambda I - \mathbf{B}(z))$ likewise has $m$ distinct roots if $z$ is close to $z_0$. Recalling \cref{thm:smooth-dependence-eigenvalues-eigenvectors} we see that near $z_0$ there exist smooth functions $\bar\lambda_1(z) < \cdots < \bar\lambda_m(z)$ and unit vectors $\{\mathbf{r}_k(z), \mathbf{l}_k(z)\}_{k=1}^m$ satisfying
$$\bar\lambda_k(z_0) = \lambda_k(z_0), \quad \mathbf{r}_k(z_0) = r_k(z_0), \quad \mathbf{l}_k(z_0) = l_k(z_0) \quad (k=1,\ldots,m),$$
and
$$(18) \qquad \begin{cases} \mathbf{B}(z)\mathbf{r}_k(z) = \bar\lambda_k(z)\mathbf{r}_k(z) \\ \mathbf{l}_k(z)\mathbf{B}(z) = \bar\lambda_k(z)\mathbf{l}_k(z) \end{cases} \quad (k=1,\ldots,m).$$
Note that $\{\mathbf{r}_k(z)\}$, $\{\mathbf{l}_k(z)\}_{k=1}^m$ are bases of $\mathbb{R}^m$, and also
$$(19) \qquad \mathbf{l}_l(z)\cdot\mathbf{r}_k(z) = 0 \quad (l\neq k).$$

3. Equation (16) will hold provided $\sigma = \bar\lambda_k(z)$ for some $k\in\{1,\ldots,m\}$, and $(z-z_0)$ is parallel to $\mathbf{r}_k(z)$. In light of (19), these conditions are equivalent to asking
$$(20) \qquad \mathbf{l}_l(z)\cdot(z-z_0) = 0 \quad (l\neq k).$$
These equalities amount to $m-1$ equations for the $m$ unknown components of $z$, which we intend to solve using the Implicit Function Theorem. So define $\Phi_k:\mathbb{R}^m\to\mathbb{R}^{m-1}$ by setting
$$\Phi_k(z) := (\ldots, \hat{\mathbf{l}}_{k-1}(z)\cdot(z-z_0), \hat{\mathbf{l}}_{k+1}(z)\cdot(z-z_0), \ldots).$$
Now $\Phi_k(z_0) = 0$ and
$$D\Phi_k(z_0) = \begin{pmatrix} \mathbf{l}_1(z_0) \\ \vdots \\ \mathbf{l}_{k-1}(z_0) \\ \mathbf{l}_{k+1}(z_0) \\ \vdots \\ \mathbf{l}_m(z_0) \end{pmatrix}_{(m-1)\times m},$$
the entries of this matrix being regarded as row vectors. Since the vectors $\{\mathbf{l}_k(z_0)\}_{k=1}^m$ form a basis of $\mathbb{R}^m$, we see
$$\operatorname{rank} D\Phi_k(z_0) = m-1.$$
Accordingly, there exists a smooth curve $\phi_k:\mathbb{R}\to\mathbb{R}^m$ such that
$$(21) \qquad \phi_k(0) = z_0$$
and
$$(22) \qquad \Phi_k(\phi_k(t)) = 0 \quad \text{for all } t \text{ close to } 0.$$
The path of the curve $\phi_k(\cdot)$ for $t$ near zero defines $S_k(z_0)$. We may reparameterize as necessary to ensure
$$(23) \qquad |\dot\phi_k(t)| = 1 \quad \left(\cdot = \frac{d}{dt}\right).$$

4. Now (20)--(22) imply
$$(24) \qquad \phi_k(t) = z_0 + \mu(t)\hat{\mathbf{r}}_k(\phi_k(t))$$
for all $t$ near zero, where $\mu:\mathbb{R}\to\mathbb{R}$ is a smooth function satisfying $\mu(0)=0$, $\dot\mu(0)=1$. Differentiating (24) with respect to $t$ and setting $t=0$, we thus find
$$(25) \qquad \dot\phi_k(0) = \mathbf{r}_k(z_0).$$
Hence the curve $S_k(z_0)$ has tangent $\mathbf{r}_k(z_0)$ at $z_0$. Assertion (i) is proved.

5. In light of the foregoing analysis, there exists a smooth function $\sigma:\mathbb{R}^m\times\mathbb{R}^m\to\mathbb{R}$ such that
$$(26) \qquad \mathbf{F}(\phi_k(t)) - \mathbf{F}(z_0) = \sigma(\phi_k(t),z_0)(\phi_k(t)-z_0)$$
for all $t$ close to zero. Differentiating with respect to $t$ and setting $t=0$, we deduce from (21) that
$$DF(z_0)\dot\phi_k(0) = \sigma(z_0,z_0)\dot\phi_k(0).$$
In light of (25), we see that $\sigma(z_0,z_0) = \lambda_k(z_0)$. This establishes assertion (ii).

6. Now write $\sigma(t) := \sigma(\phi_k(t),z_0)$; so that (26) reads
$$\mathbf{F}(\phi_k(t)) - \mathbf{F}(z_0) = \sigma(t)(\phi_k(t)-z_0).$$
Differentiate twice with respect to $t$:
$$(D^2\mathbf{F}(\phi_k(t))\dot\phi_k(t))\dot\phi_k(t) + DF(\phi_k(t))\ddot\phi_k(t) = \dot\sigma(t)(\phi_k(t)-z_0) + 2\dot\sigma(t)\dot\phi_k(t) + \sigma(t)\ddot\phi_k(t).$$
Evaluate this expression at $t=0$ and recall $\sigma(0)=\lambda_k(z_0)$, $\phi_k(0)=z_0$, $\dot\phi_k(0)=r_k(z_0)$:
$$(27) \qquad (2\dot\sigma(0)I - D^2\mathbf{F}(z_0)r_k(z_0))r_k(z_0) = (DF(z_0)-\lambda_k(z_0)I)\ddot\phi_k(0).$$

7. Let $\psi_k(t) = v(t)$ be a unit speed parametrization of the rarefaction curve $R_k(z_0)$ near $z_0$ (as in (6) above). Then
$$(28) \qquad \psi_k(0) = z_0, \quad \dot\psi_k(t) = r_k(\psi_k(t)).$$
Thus
$$DF(\psi_k(t))r_k(t) = \lambda_k(t)r_k(t),$$
for
$$\lambda_k(t) := \lambda_k(\psi_k(t)), \quad r_k(t) := r_k(\psi_k(t)).$$
Next differentiate with respect to $t$ and set $t=0$:
$$(29) \qquad (D^2\mathbf{F}(z_0)r_k(z_0) - \dot\lambda_k(0)I)r_k(z_0) = -(DF(z_0)-\lambda_k(z_0)I)\dot r_k(0).$$
Add (27) and (29), to obtain
$$(30) \qquad (2\dot\sigma(0) - \dot\lambda_k(0))r_k(z_0) = (DF(z_0)-\lambda_k(z_0)I)(\ddot\phi_k(0)-\dot r_k(0)).$$
Take the dot product with $l_k(z_0)$ and observe $l_k\cdot r_k\neq 0$, to conclude
$$(31) \qquad 2\dot\sigma(0) = \dot\lambda_k(0).$$
We deduce from (31) that
$$2\sigma(t) - \lambda_k(z_0) - \lambda_k(t) = O(t^2) \quad \text{as } t\to 0.$$
Assertion (iii) follows. $\square$
\end{proof}

We see from \cref{thm:structure-of-shock-set}(iii) that the curves $R_k(z_0)$ and $S_k(z_0)$ agree at least to first order at $z_0$. Next is the assertion that in the linearly degenerate case these curves in fact coincide.

\begin{theorem}[Linear degeneracy]
\label{thm:linear-degeneracy-shock-rarefaction-coincide}
\dcref{ch11:11.2-thm-3}
\group{Conservation Laws}
\uses{thm:structure-of-shock-set,def:rarefaction-curve,def:genuinely-nonlinear-linearly-degenerate}
Suppose for some $k\in\{1,\ldots,m\}$ that the pair $(\lambda_k, r_k)$ is linearly degenerate. Then for each $z_0\in\mathbb{R}^m$,

(i) $R_k(z_0) = S_k(z_0)$

and

(ii) $\sigma(z,z_0) = \lambda_k(z) = \lambda_k(z_0)$ for all $z\in S_k(z_0)$.
\end{theorem}

\begin{proof}
Let $\mathbf{v} = \mathbf{v}(s)$ solve the ODE
$$\begin{cases} \dot{\mathbf{v}}(s) = \mathbf{r}_k(\mathbf{v}(s)) & (s\in\mathbb{R}) \\ \mathbf{v}(0) = z_0. \end{cases}$$
Then the mapping $s\mapsto\lambda_k(\mathbf{v}(s))$ is constant, and so
$$\mathbf{F}(\mathbf{v}(s)) - \mathbf{F}(z_0) = \int_0^s DF(\mathbf{v}(t))\dot{\mathbf{v}}(t)\,dt = \int_0^s DF(\mathbf{v}(t))\mathbf{r}_k(\mathbf{v}(t))\,dt = \int_0^s \lambda_k(\mathbf{v}(t))\mathbf{r}_k(\mathbf{v}(t))\,dt = \lambda_k(z_0)\int_0^s\dot{\mathbf{v}}(t)\,dt = \lambda_k(z_0)(\mathbf{v}(s)-z_0). \qquad \square$$
\end{proof}

\subsubsection*{b. Contact discontinuities, shock waves}

We next undertake to analyze in light of \cref{thm:structure-of-shock-set,thm:linear-degeneracy-shock-rarefaction-coincide} the possibility of solving Riemann's problem by joining two given states $u_l$ and $u_r$ by some kind of shock wave.

\textbf{Contact discontinuities.} Suppose first that $(\lambda_k,r_k)$ is linearly degenerate and
$$(32) \qquad u_r \in S_k(u_l).$$

\begin{figure}[h]\centering
% [FIGURE: Diagram showing parallel diagonal lines (representing characteristics) with a horizontal line at the bottom (the discontinuity), labeled "k-contact discontinuity"]
\end{figure}

We then define an integral solution of our system of conservation laws by setting
$$\mathbf{u}(x,t) = \begin{cases} \mathbf{u}_l & \text{if } x<\sigma t \\ \mathbf{u}_r & \text{if } x>\sigma t, \end{cases}$$
for
$$\sigma = \sigma(\mathbf{u}_r,\mathbf{u}_l) = \lambda_k(\mathbf{u}_l) = \lambda_k(\mathbf{u}_r).$$
Now observe from our analysis in \cref{rem:characteristic-ode-closure} that since $\lambda_k(\mathbf{u}_l) = \lambda_k(\mathbf{u}_r) = \sigma$, the projected characteristics to the left and right are \textit{parallel} to the line of discontinuity. We interpret this situation physically by saying that fluid particles do not cross the discontinuity. The line $x=\sigma t$ is called a \textit{$k$-contact discontinuity}.

\textbf{Shock waves.} We next turn our attention to the case $(\lambda_k, r_k)$ is \textit{genuinely nonlinear}, and
$$(35) \qquad \mathbf{u}_r \in S_k(\mathbf{u}_l)$$
as before. If we consider the integral solution
$$\mathbf{u}(x,t) = \begin{cases} \mathbf{u}_l & \text{if } x<\sigma t \\ \mathbf{u}_r & \text{if } x>\sigma t, \end{cases}$$
for
$$\sigma = \sigma(\mathbf{u}_r,\mathbf{u}_l),$$
we see that there are two essentially different cases according as to whether
$$(38) \qquad \lambda_k(\mathbf{u}_r) < \lambda_k(\mathbf{u}_l)$$
or else
$$(39) \qquad \lambda_k(u_l) < \lambda_k(u_r).$$
Now in view of assertion (iii) from \cref{thm:structure-of-shock-set}, we have then either
$$(40) \qquad \lambda_k(u_r) < \sigma(u_r,u_l) < \lambda_k(u_l)$$
or
$$(41) \qquad \lambda_k(u_l) < \sigma(u_r,u_l) < \lambda_k(u_r),$$
provided $u_r$ is close enough to $u_l$.

This dichotomy is reminiscent of a corresponding situation in \S3.4, for a scalar conservation law. By analogy with the entropy conditions introduced there, let us hereafter agree to reject the inequalities (41) as allowing for ``nonphysical shocks'' from which characteristics emanate as we move forward in time. We rather take (40) as being physically correct. The informal viewpoint is that then the characteristics from the left and right run into the line of discontinuity, whereupon ``information is lost'' and so ``entropy increases''. This interpretation was largely justified mathematically in \S3.4 with our uniqueness theorem (\cref{thm:uniqueness-entropy-solutions}) for weak solutions that satisfied this sort of entropy condition.

Refocusing our attention again to systems, we therefore agree to regard (40) as the correct inequalities to be satisfied:

\begin{definition}[Admissible pair, Lax entropy condition]
\label{def:admissible-pair-lax-entropy-condition}
\dcref{ch11:11.2-def-4}
\group{Conservation Laws}
\uses{def:shock-set,thm:structure-of-shock-set}
Assume the pair $(\lambda_k, r_k)$ is genuinely nonlinear at $u_l$. We say that the pair $(u_l, u_r)$ is \textit{admissible} provided
$$u_r \in S_k(u_l)$$
and
$$(42) \qquad \lambda_k(u_r) < \sigma(u_r,u_l) < \lambda_k(u_l).$$
We refer to (42) as the \textit{Lax entropy condition}. If $(u_l,u_r)$ is admissible, we call our solution $\mathbf{u}$ defined above a \textit{$k$-shock wave}.
\end{definition}

By analogy with our decomposition of $R_k(z_0)$ into $R_k^\pm(z_0)$, let us introduce this

\begin{definition}[$S_k^+$ and $S_k^-$]
\label{def:shock-set-plus-minus}
\dcref{ch11:11.2-def-5}
\group{Conservation Laws}
\uses{def:shock-set,def:admissible-pair-lax-entropy-condition}
If the pair $(\lambda_k, r_k)$ is genuinely nonlinear, we write
$$S_k^+(z_0) := \{z\in S_k(z_0) \mid \lambda_k(z_0) < \sigma(z,z_0) < \lambda_k(z)\}$$

\begin{figure}[h]\centering
% [FIGURE: A shock curve diagram showing a wavy curve with three marked regions: $S_k^-(z_0)$ on the left branch, a point $z_0$ at the center, and $S_k^+(z_0)$ on the right branch]
\end{figure}

and
$$S_k^-(z_0) := \{z\in S_k(z_0) \mid \lambda_k(z) < \sigma(z,z_0) < \lambda_k(z_0)\}.$$
Then
$$S_k(z_0) = S_k^+(z_0) \cup \{z_0\} \cup S_k^-(z_0)$$
near $z_0$. Note then that the pair $(u_l, u_r)$ is admissible if and only if
$$u_r \in S_k^-(u_l).$$
\end{definition}

\subsection{Local solution of Riemann's problem}\label{sec:local-solution-riemanns-problem}

Next we glue together the physically relevant parts of the rarefaction and shock curves.

\begin{definition}[The curve $T_k$]
\label{def:t-curve}
\dcref{ch11:11.2-def-6}
\group{Conservation Laws}
\uses{def:rarefaction-curve,def:shock-set-plus-minus,def:genuinely-nonlinear-linearly-degenerate}
(i) If the pair $(\lambda_k, r_k)$ is genuinely nonlinear, write
$$T_k(z_0) := R_k^+(z_0) \cup \{z_0\} \cup S_k^-(z_0).$$

(ii) If the pair $(\lambda_k, r_k)$ is linearly degenerate, we set
$$T_k(z_0) := R_k(z_0) = S_k(z_0).$$
\end{definition}

Owing to \cref{thm:structure-of-shock-set}(ii) the curve $T_k(z_0)$ is $C^1$. Employing this notation we see that nearby states $u_l$ and $u_r$ can be joined by a $k$-rarefaction wave, a shock wave or a contact discontinuity provided

$$(43) \qquad u_l \in T_k(u_r).$$

We now at last ask if we can find a solution to Riemann's problem provided only that $u_r$ is close to $u_l$ (but (43) may fail for each $k=1,\ldots,m$). The hope is that by moving along various paths $T_k$ for different values of $k$, we may be able to connect $u_l$ to $u_r$, utilizing a sequence of rarefaction waves, shock waves, and/or contact discontinuities.

\begin{theorem}[Local solution of Riemann's problem]
\label{thm:local-solution-riemann-problem}
\dcref{ch11:11.2-thm-4}
\group{Conservation Laws}
\uses{def:t-curve,thm:existence-k-rarefaction-waves,def:admissible-pair-lax-entropy-condition,def:eigenvalue-eigenvector-notation-hyperbolic}
Assume for each $k=1,\ldots,m$ that the pair $(\lambda_k, r_k)$ is either genuinely nonlinear or else linearly degenerate. Suppose further the left state $u_l$ is given. Then for each right state $u_r$ sufficiently close to $u_l$ there exists an integral solution $\mathbf{u}$ of Riemann's problem, which is constant on lines through the origin.
\end{theorem}

\begin{figure}[h]\centering
% [FIGURE: A curve with three key points labeled: $s_k^-(z_0)$ on the left portion of the curve, $z_0$ at the center, and $R_k^+(z_0)$ on the right portion. The curve is smooth and shows arrows at both ends indicating the direction along the curve.]
\end{figure}

Structure of the T-curve

\begin{proof}
1. We intend to apply the Inverse Function Theorem to a mapping $\Phi:\mathbb{R}^m\to\mathbb{R}^m$, defined near $0$ as follows.

First, for each family of curves $T_k$ $(k=1,\ldots,m)$ choose the nonsingular parameter $\tau_k$ to measure arc length; that is, if $z,\tilde z\in\mathbb{R}^m$ with $\tilde z\in T_k(z)$, then
$$\tau_k(\tilde z) - \tau_k(z) = \text{(signed) distance from } z \text{ to } \tilde z \text{ along the curve } T_k(z).$$
We take the plus sign for $\tau_k(\tilde z)$ if $\tilde z\in R_k^+(z)$, the minus sign if $\tilde z\in S_k^-(z)$.

2. Given then $t=(t_1,\ldots,t_m)\in\mathbb{R}^m$, with $|t|$ small, we define $\Phi(t)=z$ as follows. First, temporarily write
$$(44) \qquad u_l = z_0.$$
Then choose states $z_1,\ldots,z_m$ to satisfy
$$(45) \qquad \begin{cases} z_1\in T_1(z_0), & \tau_1(z_1)-\tau_1(z_0)=t_1, \\ z_2\in T_2(z_1), & \tau_2(z_2)-\tau_2(z_1)=t_2, \\ & \vdots \\ z_m\in T_m(z_{m-1}), & \tau_m(z_m)-\tau_m(z_{m-1})=t_m. \end{cases}$$
Now write
$$(46) \qquad z=z_m,$$
and define
$$(47) \qquad \Phi(t) = z.$$
Note $\Phi$ is $C^1$, and
$$(48) \qquad \Phi(0) = z_0.$$

3. We claim
$$(49) \qquad D\Phi(0) \text{ is nonsingular.}$$
To see this, observe
$$\Phi(0,\ldots,t_k,\ldots,0) - \Phi(0,\ldots,0) = t_k r_k(z_0) + o(t_k) \text{ as } t_k\to 0.$$
Thus
$$\frac{\partial\Phi}{\partial t_k}(0) = r_k(z_0) \quad (k=1,\ldots,m),$$
and so
$$D\Phi(0) = (\mathbf{r}_1(z_0),\ldots,\mathbf{r}_m(z_0))_{m\times m},$$
the entries regarded as column vectors. This matrix is nonsingular, since $\{\mathbf{r}_k(z_0)\}_{k=1}^m$ is a basis.

4. In light of (49), the Inverse Function Theorem applies: for each state $u_r$ sufficiently close to $u_l$ there exists a unique parameter $t=(t_1,\ldots,t_m)$ close to zero such that $\Phi(t)=u_r$.

Recall next that if $z_{k-1}$ and $z_k$ are joined by a $k$-rarefaction wave, this wave is
$$\begin{cases} z_{k-1} & \text{if } \frac{x}{t}<\lambda_k(z_{k-1}) \\ G_k\left(\frac{x}{t}\right) & \text{if } \lambda_k(z_{k-1})<\frac{x}{t}<\lambda_k(z_k), \quad \text{for } G_k=(F_k')^{-1} \\ z_k & \text{if } \lambda_k(z_k)<\frac{x}{t}. \end{cases}$$
Moreover if $z_{k-1},z_k$ are joined by a $k$-shock, it has the form
$$\begin{cases} z_{k-1} & \text{if } \frac{x}{t}<\sigma(z_k,z_{k-1}) \\ z_k & \text{if } \sigma(z_k,z_{k-1})<\frac{x}{t}, \end{cases}$$
where $\lambda_k(z_k)<\sigma(z_k,z_{k-1})<\lambda_k(z_{k-1})$. In both cases the waves are constant outside the regions $\lambda_k(z_0)-\epsilon<\frac{x}{t}<\lambda_k(z_0)+\epsilon$, for small $\epsilon>0$, provided $z_k,z_{k-1}$ are close enough to $z_0$. This is true for $k=1,\ldots,m$.

Since $\lambda_1(z_0)<\cdots<\lambda_m(z_0)$, we see then that the rarefactions, shock waves and/or contact discontinuities connecting $u_l=z_0$ to $z_1$, $z_1$ to $z_2$, $z_2$ to $z_3$, \ldots, $z_{m-1}$ to $z_m=u_r$ do not intersect. $\square$
\end{proof}

\section{Systems of Two Conservation Laws}\label{sec:systems-two-conservation-laws}

In this section we more deeply analyze the initial-value problem for $m=2$, which is to say, for a pair of conservation laws:

$$(1) \qquad \begin{cases} u^1_t + F^1(u^1,u^2)_x = 0 & \text{in } \mathbb{R}\times(0,\infty) \\ u^2_t + F^2(u^1,u^2)_x = 0 \\ u^1=g^1, u^2=g^2 & \text{on } \mathbb{R}\times\{t=0\}. \end{cases}$$

Here $\mathbf{F}=(F^1,F^2)$, $\mathbf{g}=(g^1,g^2)$, $\mathbf{u}=(u^1,u^2)$.

\subsection{Riemann invariants}\label{sec:riemann-invariants}

Our intention is first to demonstrate that we can transform (1) into a much simpler form by performing an appropriate nonlinear change of dependent variables. The idea is to find two functions $w^1,w^2:\mathbb{R}^2\to\mathbb{R}$ with nice properties along the rarefaction curves $R_1,R_2$:

\begin{definition}[$i$-th Riemann invariant]
\label{def:riemann-invariant}
\dcref{ch11:11.3-def-1}
\group{Conservation Laws}
\uses{def:eigenvalue-eigenvector-notation-hyperbolic}
We say
$$w^i:\mathbb{R}^2\to\mathbb{R}$$
is an $i^{\text{th}}$-\textit{Riemann invariant} provided
$$(2) \qquad Dw^i(z) \text{ is parallel to } l_j(z) \quad (z\in\mathbb{R}^2, i\neq j).$$
\end{definition}

We will see momentarily how useful condition (2) is, but let us first pause to ask whether Riemann invariants exist. It turns out that since we are now taking $m=2$, this is easy. Indeed, because $l_j(z)\cdot r_i(z)=0$ $(i\neq j)$, we see (2) is equivalent in $\mathbb{R}^2$ to the statement

$$(2') \qquad Dw^i(z)\cdot r_i(z) = 0 \quad (i=1,2, z\in\mathbb{R}^2);$$

which is to say

$$(3) \qquad w^i \text{ is constant along the rarefaction curve } R_i \quad (i=1,2).$$

In particular, any smooth function $w^i$ satisfying (3) satisfies also (2'), (2) and so is an $i^{\text{th}}$-Riemann invariant.

\begin{remark}[Riemann invariants for $m>2$]
\label{rem:riemann-invariants-mgt2-remark}
\dcref{ch11:11.3-rem-1}
\group{Conservation Laws}
\uses{def:riemann-invariant}
In the case that $m>2$, Riemann invariants do not in general exist. $\square$
\end{remark}

Now we can regard $w=(w_1,w_2) = (w^1(z_1,z_2), w^2(z_1,z_2))$ as being new coordinates on the state space $\mathbb{R}^2$, replacing $z=(z_1,z_2)$. More precisely, we define $\mathbf{w}:\mathbb{R}^2\to\mathbb{R}^2$ by setting

$$(4) \qquad \mathbf{w}(z) = \mathbf{w}(z_1,z_2) = (w^1(z_1,z_2), w^2(z_1,z_2)).$$

The inverse mapping is $\mathbf{z}(w) = \mathbf{z}(w_1,w_2) = (z^1(w_1,w_2), z^2(w_1,w_2))$.

Let us now utilize the transformation (4) to simplify our system of two conservation laws (1). For this, let us suppose henceforth $\mathbf{u} = (u^1,u^2)$ is a smooth solution of (1). We now change dependent variables by setting

$$(5) \qquad \mathbf{v}(x,t) := \mathbf{w}(\mathbf{u}(x,t)) \quad (x\in\mathbb{R}, t>0).$$

What system of PDE does $\mathbf{v} = (v^1,v^2)$ satisfy?

\begin{theorem}[Conservation laws and Riemann invariants]
\label{thm:conservation-laws-riemann-invariants}
\dcref{ch11:11.3-thm-1}
\group{Conservation Laws}
\uses{def:riemann-invariant}
The functions $v^1, v^2$ solve the system
$$(6) \qquad \begin{cases} v^1_t + \lambda_2(\mathbf{u})v^1_x = 0 \\ v^2_t + \lambda_1(\mathbf{u})v^2_x = 0 \end{cases} \quad \text{in } \mathbb{R}\times(0,\infty).$$
The point is that the system (6), although not in conservation law form, is in many ways rather simpler than (1). Note in particular that whereas the PDE for $u^1$ involves the term $u^2_x$, the PDE for $v^1$ does \textit{not} entail $v^2_x$. Similarly, the PDE for $v^2$ does not involve $v^1_x$.
\end{theorem}

\begin{proof}
According to (5), we see that for $i=1,2$, $i\neq j$,
$$v^i_t + \lambda_j(\mathbf{u})v^i_x = Dw^i(\mathbf{u})\cdot\mathbf{u}_t + \lambda_j(\mathbf{u})Dw^i(\mathbf{u})\cdot\mathbf{u}_x = Dw^i(\mathbf{u})\cdot(-F(\mathbf{u})_x + \lambda_j(\mathbf{u})\mathbf{u}_x) = Dw^i(\mathbf{u})\cdot(-DF(\mathbf{u}) + \lambda_j(f)I)\mathbf{u}_x = 0,$$
since, by definition, $Dw^i$ is parallel to $\mathbf{l}_j$. $\square$
\end{proof}

\begin{remark}[Riemann invariants along characteristics]
\label{rem:riemann-invariant-characteristic-curves}
\dcref{ch11:11.3-rem-2}
\group{Conservation Laws}
\uses{thm:conservation-laws-riemann-invariants,def:genuinely-nonlinear-linearly-degenerate}
(i) We can interpret the system of PDE (6) by introducing the ODE
$$(7) \qquad \dot x_i(s) = \lambda_j(\mathbf{u}(x_i(s),s)) \quad (s\geq 0)$$
for $i=1,2$, $j\neq i$. Then we see from (6) that
$$(8) \qquad v^i \text{ is constant along the curve } (x_i(s),s) \quad (s\geq 0)$$
for $i=1,2$.

(ii) Recall from \S11.2 our condition of genuine nonlinearity reads
$$(9) \qquad D\lambda_i(z)\cdot\mathbf{r}_i(z) \neq 0 \quad (z\in\mathbb{R}^2, i=1,2).$$
Since we can also regard $\lambda_i$ as a function of $w=(w_1,w_2)$, we can rewrite (9) to read
$$(10) \qquad \frac{\partial\lambda_i}{\partial w_j} \neq 0 \quad (w\in\mathbb{R}^2, i\neq j).$$
To see (9) is equivalent to (10), observe that if (10) fails, then
$$(11) \qquad 0 = \frac{\partial\lambda_i}{\partial w_j} = \sum_{k=1}^2 \frac{\partial\lambda_i}{\partial z_k}\frac{\partial z^k}{\partial w_j}.$$
But since $\sum_{k=1}^2 \frac{\partial w^k}{\partial z_k}\frac{\partial z^k}{\partial w_j} = \delta_{ij} = 0$ for $i\neq j$, we see that (11) asserts that $D\lambda_i$ is parallel to $Dw^i$. However, $Dw^i$ is perpendicular to $\mathbf{r}_i$, and so we obtain a contradiction to (9). Hence (9) implies (10), and the converse implication is established in the same way. $\square$
\end{remark}

\begin{example}[Barotropic compressible gas dynamics]
\label{ex:barotropic-gas-dynamics-riemann-invariants}
\dcref{ch11:11.3-ex-1}
\group{Conservation Laws}
\uses{def:riemann-invariant,ex:euler-equations-1d-system,rem:riemann-invariant-characteristic-curves}
We illustrate the foregoing ideas by examining in detail Euler's equations for barotropic compressible gas dynamics, a special case of the general Euler equations (\cref{ex:euler-equations-1d-system}) when the internal energy $e$ is constant. The relevant PDE are
$$(12) \qquad \begin{cases} \rho_t + (\rho v)_x = 0 & \text{(conservation of mass)} \\ (\rho v)_t + (\rho v^2+p)_x = 0 & \text{(conservation of momentum)}, \end{cases}$$
where we now assume
$$(13) \qquad p = p(\rho)$$
for some smooth function $p:\mathbb{R}\to\mathbb{R}$. Formula (13) is called a \textit{barotropic equation of state}. We assume the strict hyperbolicity condition
$$(14) \qquad p' > 0.$$
Setting $\mathbf{u}=(u^1,u^2)=(\rho,\rho v)$, we can rewrite (12), (13) to read
$$\mathbf{u}_t + \mathbf{F}(\mathbf{u})_x = 0,$$
for
$$\mathbf{F} = (F^1,F^2) = (z_2, (z_2)^2/z_1 + p(z_1))$$
and $z=(z_1,z_2)$, provided $z_1>0$. Then
$$D\mathbf{F} = \begin{pmatrix} 0 & 1 \\ -(z_2/z_1)^2 + p'(z_1) & 2z_2/z_1 \end{pmatrix}.$$
Consequently
$$(15) \qquad \lambda_1 = \frac{z_2}{z_1} - p'(z_1)^{1/2}, \quad \lambda_2 = \frac{z_2}{z_1} + p'(z_1)^{1/2}.$$
In physical notation:
$$(16) \qquad \lambda_1 = v-\sigma, \quad \lambda_2 = v+\sigma,$$
for the sound speed
$$(17) \qquad \sigma := p'(\rho)^{1/2}.$$
Remembering (7), we consider next the ODE
$$(18) \qquad \dot x_1(t) = v(x_1(t),t) + \sigma(x_1(t),t),$$
$$(19) \qquad \dot x_2(t) = v(x_2(t),t) - \sigma(x_2(t),t),$$
where $\sigma(x,t) := p'(\rho(x,t))^{\frac12}$, $t\geq 0$. We know from (8) that the Riemann invariant $w^1=w^1(\mathbf{u})$ is constant along the trajectories of (18) and $w^2=w^2(\mathbf{u})$ is constant along trajectories of (19).

To compute $w^1$ and $w^2$ directly, let us carry out some computations on the system (12). First we transform (12) in nondivergence form:
$$(20) \qquad \rho_t + \rho v_x + \rho_x v = 0,$$
$$(21) \qquad \rho_t v + \rho v_t + \rho_x v^2 + 2\rho v v_x + p_x = 0.$$
Multiplying (20) by $\sigma^2 = p'(\rho)$ and recalling (13) gives us
$$(22) \qquad p_t + vp_x + \sigma^2\rho v_x = 0.$$
In addition (20), (21) combine to yield
$$(23) \qquad \rho v_t + \rho v v_x + p_x = 0.$$
We now manipulate (22), (23) so that the directions $\lambda_1,\lambda_2 = v\mp\sigma$ appear explicitly. To accomplish this, we multiply (23) by $\sigma$ and then add to and subtract from (22):
$$(24) \qquad \begin{cases} p_t + (v+\sigma)p_x + \rho\sigma(v_t+(v+\sigma)v_x) = 0 \\ p_t + (v-\sigma)p_x - \rho\sigma(v_t+(v-\sigma)v_x) = 0. \end{cases}$$
We then deduce from (24) that
$$(25) \qquad \begin{cases} \frac{d}{dt}[p(x_1(t),t)] + \rho(x_1(t),t)\sigma(x_1(t),t)\frac{d}{dt}[v(x_1(t),t)] = 0 \\ \frac{d}{dt}[p(x_2(t),t)] - \rho(x_2(t),t)\sigma(x_2(t),t)\frac{d}{dt}[v(x_2(t),t)] = 0. \end{cases}$$
As $\frac{dp}{dt} = \sigma^2\frac{d\rho}{dt}$, we see
$$(26) \qquad \frac{\sigma}{\rho}\frac{d\sigma}{dt} = \frac{dv}{dt} = 0 \quad \text{along the trajectories of (18), (19),}$$
provided $\rho>0$.

Think now of the Riemann invariants as functions of $\rho$ and $v$. Then since $w^1=w^1(\rho,v)$ is constant along the curve determined by $x_1(\cdot)$, we have
$$0 = \frac{d}{dt}[w^1(\rho(x_1(t),t),v(x_1(t),t))] = \frac{\partial w^1}{\partial\rho}\frac{d}{dt}[\rho(x_1(t),t)] + \frac{\partial w^1}{\partial v}\frac{d}{dt}[v(x_1(t),t)].$$
This is consistent with (26) if
$$\frac{\partial w^1}{\partial\rho} = \frac{\sigma(\rho)}{\rho}, \quad \frac{\partial w^1}{\partial v} = 1.$$
We similarly deduce
$$\frac{\partial w^2}{\partial\rho} = \frac{\sigma(\rho)}{\rho}, \quad \frac{\partial w^2}{\partial v} = -1.$$
Integrating, we conclude that the Riemann invariants are, up to additive constants,
$$w^1 = \int_1^\rho \frac{\sigma(s)}{s}\,ds + v, \quad w^2 = \int_1^\rho \frac{\sigma(s)}{s}\,ds - v.$$
We leave it as an exercise to check that $w^1, w^2$, taken now as functions of $z=(z_1,z_2)$, satisfy the definition of Riemann invariants. $\square$
\end{example}

\subsection{Nonexistence of smooth solutions}\label{sec:nonexistence-smooth-solutions-systems}

Illustrating now the usefulness of Riemann invariants, we establish the following criterion for the nonexistence of a smooth solution:

\begin{theorem}[Riemann invariants and blow-up]
\label{thm:riemann-invariants-blowup}
\dcref{ch11:11.3-thm-2}
\group{Conservation Laws}
\uses{def:riemann-invariant,thm:conservation-laws-riemann-invariants}
Assume $g$ is smooth, with compact support. Suppose also the genuine nonlinearity condition
$$(27) \qquad \frac{\partial\lambda_i}{\partial w_j} > 0 \text{ in } \mathbb{R}^2 \quad (i=1,2,\ i\neq j)$$
holds. Then the initial-value problem (1) cannot have a smooth solution $u$ existing for all times $t\geq 0$ if
$$(28) \qquad \text{either } v^1_x < 0 \text{ or } v^2_x < 0 \text{ somewhere on } \mathbb{R}\times\{t=0\}.$$
\end{theorem}

\begin{proof}
1. Assume for the time being that $\mathbf{u}$ is a smooth solution of (1). Write
$$(29) \qquad a := v^1_x, \quad b := v^2_x,$$
where $\mathbf{v} = \mathbf{w}(\mathbf{u})$, $\mathbf{v}=(v^1,v^2)$, solves the system of PDE (6). We differentiate the first equation of (6) with respect to $x$, to compute:
$$(30) \qquad a_t + \lambda_2 a_x + \frac{\partial\lambda_2}{\partial w_1}a^2 + \frac{\partial\lambda_2}{\partial w_2}ab = 0.$$
We employ then the second equation of (6), which we rewrite as
$$v^2_t + \lambda_2 v^2_x = (\lambda_2-\lambda_1)b.$$
Substituting this expression into (30) gives
$$(31) \qquad a_t + \lambda_2 a_x + \frac{\partial\lambda_2}{\partial w_1}a^2 + \left[\frac{1}{\lambda_2-\lambda_1}\frac{\partial\lambda_2}{\partial w_2}(v^2_t+\lambda_2 v^2_x)\right]a = 0.$$

2. To integrate (31), fix $x_0\in\mathbb{R}$ and set
$$(32) \qquad \xi(t) := \exp\left(\int_0^t \frac{1}{\lambda_2-\lambda_1}\frac{\partial\lambda_2}{\partial w_2}(v^2_t+\lambda_2 v^2_x)(x_1(s),s)\,ds\right),$$
where
$$(33) \qquad \begin{cases} \dot x_1(s) = \lambda_2(\mathbf{u}(x_1(s),s)) & (s\geq 0) \\ x_1(0) = x_0. \end{cases}$$
Next is the key observation from (8) that $v^1$ is constant along the curve $(x_1(s),s)$. So write
$$v^1(x_1(s),s) = v^1_0 = v^1(x_0,0) \quad (s\geq 0).$$
Thus we see that the expression $\left(\frac{1}{\lambda_2-\lambda_1}\frac{\partial\lambda_2}{\partial w_2}\right)$, considered as a function of $\mathbf{v}=\mathbf{w}(\mathbf{u})$, depends only on $v^2$. Let us set
$$\gamma(\mu) := \int_0^\mu \left(\frac{1}{\lambda_2-\lambda_1}\frac{\partial\lambda_2}{\partial w_2}\right)(v^1_0,v)\,dv.$$
Then (32), (33) imply
$$(34) \qquad \xi(t) = \exp\left(\int_0^t \frac{d}{ds}[\gamma(v^2(x_1(s),s))]\,ds\right) = \exp(\gamma(v^2(x_1(t),t)) - \gamma(v^2(x_0,0))).$$

3. We now transform (31) to read
$$\frac{d}{dt}((\xi\alpha)^{-1}) = \frac{-1}{(\xi\alpha)^2}\frac{d}{dt}(\xi\alpha) = \frac{\partial\lambda_2}{\partial w_1}\xi^{-1}$$
where $\alpha(t) := a(x_1(t),t)$ and we assume $\alpha\neq 0$. Consequently
$$(\xi(t)\alpha(t))^{-1} = (\alpha(0))^{-1} + \int_0^t \frac{\partial\lambda_2}{\partial w_1}\xi^{-1}(s)\,ds.$$
This equality in turn rearranges to become
$$(35) \qquad \alpha(t) = \alpha(0)\xi^{-1}(t)\left(1 + \alpha(0)\int_0^t \frac{\partial\lambda_2}{\partial w_1}\xi^{-1}(s)\,ds\right)^{-1}.$$

4. Now in view of the system of PDE (6), $\mathbf{v}$ is bounded. Thus we deduce from (34) that $0<\theta\leq\xi(t)\leq\Theta$ for all times $t>0$, for appropriate constants $\theta,\Theta$. Therefore it follows from (27) and (35) that $\alpha$ is bounded for all $t>0$ if and only if $\alpha(0)\geq 0$, that is, if
$$v^1_x(x_0,0) \geq 0.$$
A similar calculation holds with $v^2$ replacing $v^1$. We conclude that if either $v^1_x<0$ or else $v^2_x<0$ somewhere on $\mathbb{R}\times\{t=0\}$, there cannot then exist a smooth solution of (1), lasting for all times $t\geq 0$. $\square$
\end{proof}

\section{Entropy Criteria}\label{sec:entropy-criteria-systems}

In our study of Riemann's problem in \S11.2 we have taken Lax's entropy condition
$$(1) \qquad \lambda_k(u_r) < \sigma(u_r,u_l) < \lambda_k(u_l)$$
for some $k\in\{1,\ldots,m\}$ as the selection criteria for admissible shock waves.

There is great ongoing interest in discovering other mathematically correct and physically appropriate entropy conditions of various sorts, with the aim of applying these to more complicated integral solutions of our system of conservation laws, so as to obtain uniqueness criteria, more information concerning allowable discontinuities, etc.

One general principle, instances of which we have already seen for scalar conservation laws in \S4.5.2 and for Hamilton--Jacobi equations in \S10.1, is that physically and mathematically correct solutions should arise as the limit of solutions to the regularized system
$$(2) \qquad \mathbf{u}^\varepsilon_t + \mathbf{F}(\mathbf{u}^\varepsilon)_x - \varepsilon\mathbf{u}^\varepsilon_{xx} = 0 \quad \text{in } \mathbb{R}\times(0,\infty)$$
as $\varepsilon\to 0$. The idea is to interpret the term ``$\varepsilon\mathbf{u}^\varepsilon_{xx}$'' as providing a small viscosity effect, which will presumably ``smear out'' sharp shocks. The hope is to study various aspects of the problem (2) in the limit $\varepsilon\to 0$, and thereby to discover more general entropy criteria, to augment Lax's condition (1).

The next sections discuss aspects of this general program.

\subsection{Vanishing viscosity, traveling waves}\label{sec:vanishing-viscosity-traveling-waves}

We begin our investigation of the parabolic system (2) by first seeking a traveling wave solution, having the form

$$(3) \qquad u^\varepsilon(x,t) = v\left(\frac{x-\sigma t}{\varepsilon}\right) \quad (x\in\mathbb{R}, t\geq 0),$$

where as usual the speed $\sigma$ and profile $v$ must be found. Substituting (3) into (2), we find $v:\mathbb{R}\to\mathbb{R}^m$, $v=v(s)$, must solve the ODE

$$(4) \qquad \dot v = -\sigma v + DF(v)\dot v \quad \left(\dot{} = \frac{d}{ds}\right).$$

Assume now $u_l, u_r\in\mathbb{R}^m$ are given, and furthermore

$$(5) \qquad \lim_{s\to-\infty} v = u_l, \quad \lim_{s\to+\infty} v = u_r, \quad \lim_{s\to\pm\infty} \dot v = 0.$$

Then from (3) we deduce

$$(6) \qquad \lim_{\varepsilon\to 0} u^\varepsilon(x,t) = \begin{cases} u_l & \text{if } x<\sigma t \\ u_r & \text{if } x>\sigma t. \end{cases}$$

Hence the limit as $\varepsilon\to 0$ of our solution to (2) gives us a shock wave connecting the states $u_l, u_r$. The plan now is to study carefully the form of $\sigma$ and $v$, and thereby glean more detailed information about the structure of the shock determined by (6).

The first and primary question is whether there in fact exist $\sigma$ and $v$ solving (4), (5). Integrating (4), we deduce

$$(7) \qquad \dot v = \mathbf{F}(v) - \sigma v + c$$

for some constant $c\in\mathbb{R}^m$. We conclude from (5) that

$$(8) \qquad \mathbf{F}(u_l) - \sigma u_l + c = \mathbf{F}(u_r) - \sigma u_r + c.$$

Hence

$$(9) \qquad \mathbf{F}(u_l) - \mathbf{F}(u_r) = \sigma(u_l - u_r).$$

In view of (5), (8) and (9), our ODE (7) becomes

$$(10) \qquad \dot v = \mathbf{F}(v) - \mathbf{F}(u_l) - \sigma(v-u_l).$$

Now think of the left state $u_l$ as being given, and suppose we are trying to build a traveling wave connecting $u_l$ to a nearby state $u_r$. From (9) we see that necessarily $u_r\in S_k(u_l)$ for some $k\in\{1,\ldots,m\}$, and

$$(11) \qquad \sigma = \sigma(u_r,u_l).$$

We refine this observation as follows:

\begin{theorem}[Existence of traveling waves for genuinely nonlinear systems]
\label{thm:existence-traveling-waves-genuinely-nonlinear}
\dcref{ch11:11.4-thm-1}
\group{Conservation Laws}
\uses{def:genuinely-nonlinear-linearly-degenerate,def:shock-set,thm:structure-of-shock-set}
Assume the pair $(\lambda_k, \mathbf{r}_k)$ is genuinely nonlinear for $k=1,\ldots,m$. Let $u_r$ be selected sufficiently close to $u_l$. Then there exists a traveling wave solution of (2) connecting $u_l$ to $u_r$ if and only if
$$(12) \qquad u_r \in S_k(u_l)$$
for some $k\in\{1,\ldots,m\}$.
\end{theorem}

\begin{proof}
1. Assume first $\sigma$ and $\mathbf{v}$ solve (4), (5). Then, as noted above, necessarily $u_r\in S_k(u_l)$ for some $k\in\{1,\ldots,m\}$, $\sigma=\sigma(u_r,u_l)$. Now set
$$(13) \qquad \mathbf{G}(z) := \mathbf{F}(z) - \mathbf{F}(u_l) - \sigma(z-u_l).$$
Our ODE (10) then reads
$$(14) \qquad \dot{\mathbf{v}} = \mathbf{G}(\mathbf{v}),$$
and we have
$$(15) \qquad \mathbf{G}(u_l) = \mathbf{G}(u_r) = 0,$$
according to (9). We compute
$$D\mathbf{G}(u_l) = D\mathbf{F}(u_l) - \sigma I;$$
and so the eigenvalues of $D\mathbf{G}$ at $u_l$ are $\{\lambda_k(u_l)-\sigma\}_{k=1}^m$, with corresponding right and left eigenvectors $\{\mathbf{r}_k, l_k\}_{k=1}^m$, $\mathbf{r}_k=\mathbf{r}_k(u_l)$, $l_k=l_k(u_l)$.

2. Now since $u_r\in S_k(u_l)$ and $|u_r-u_l|$ is small, we know from \cref{thm:structure-of-shock-set}(iii) that
$$\sigma = \frac{\lambda_k(u_r)+\lambda_k(u_l)}{2} + o(|u_r-u_l|).$$
Thus
$$\lambda_k(u_l) - \sigma = \frac{\lambda_k(u_l)-\lambda_k(u_r)}{2} + o(|u_r-u_l|).$$
In order that there be an orbit of the ODE (14) connecting $u_l$ at $s=-\infty$ to $u_r\in S_k(u_l)$ at $s=+\infty$, it must be that $\lambda_k(u_l)-\sigma>0$; for otherwise the trajectory would not converge to $u_l$ as $s\to-\infty$. Thus if $|u_r-u_l|$ is small enough, $\lambda_k(u_r)<\lambda_k(u_l)$, which is to say $u_r\in S_k^-(u_l)$.

3. We omit proof of the sufficiency of condition (12): see Majda--Pego \cite{MajdaPego1985}. $\square$
\end{proof}

The preceding result employs the genuine nonlinearity assumption, but the assertion holds in general, provided we introduce an appropriate variant of Lax's entropy condition (1). So let us suppose now $u_r\in S_k(u_l)$ for some $k\in\{1,\ldots,m\}$ and furthermore

\begin{definition}[Lax's entropy criterion, general form]
\label{def:lax-entropy-criterion-general}
\dcref{ch11:11.4-def-1}
\group{Conservation Laws}
\uses{def:shock-set}
$$(16) \qquad \begin{cases} \sigma(z,u_l) > \sigma(u_r,u_l) \text{ for each } z \text{ lying} \\ \text{on the curve } S_k(u_l) \text{ between } u_r \text{ and } u_l. \end{cases}$$
Condition (16) is \textit{Lax's entropy criterion}. (Observe this condition is automatic provided $(\lambda_k, r_k)$ is genuinely nonlinear, $u_r\in S_k(u_l)$, and $u_r$ is sufficiently close to $u_l$.)
\end{definition}

We can motivate (16) by again seeking traveling wave solutions of system (2). So assume $u_l$ is given. Then provided $|u_r-u_l|$ is small enough, it turns out that there exists a traveling wave solution $u^\varepsilon(x,t) = v\left(\frac{x-\sigma t}{\varepsilon}\right)$, $v$ solving (4), (5), if and only if the entropy condition (16) is satisfied. See Conlon \cite{Conlon1980} for a proof.

To make this all a bit clearer, we next present in detail a specific application.

\begin{example}[Traveling waves for the $p$-system]
\label{ex:traveling-waves-p-system}
\dcref{ch11:11.4-ex-1}
\group{Conservation Laws}
\uses{ex:p-system-definition,thm:existence-traveling-waves-genuinely-nonlinear,def:lax-entropy-criterion-general}
Let us consider again the $p$-system
$$(17) \qquad \begin{cases} u^1_t - u^2_x = 0 & \text{(compatibility condition)} \\ u^2_t - p(u^1)_x = 0 & \text{(Newton's law)}, \end{cases}$$
under the usual strict hyperbolicity condition
$$(18) \qquad p' > 0.$$
We investigate the existence of traveling wave solutions to the regularized system
$$(19) \qquad \begin{cases} u^{\varepsilon,1}_t - u^{\varepsilon,2}_x = 0 \\ u^{\varepsilon,2}_t - p(u^{\varepsilon,1})_x = \varepsilon u^{\varepsilon,2}_{xx}. \end{cases}$$
Notice we have added the viscosity term only to the second equation. This makes sense physically, as the first line of (17) is only a mathematical compatibility condition.

Assume now $u^\varepsilon = v\left(\frac{x-\sigma t}{\varepsilon}\right)$ is a traveling wave solution of (19), with
$$(20) \qquad \lim_{s\to-\infty} v = u_l, \quad \lim_{s\to\infty} v = u_r, \quad \lim_{s\to\pm\infty} \dot v = 0.$$
Writing $\mathbf{v} = (v^1,v^2)$, we compute from (19) that
$$(21) \qquad \begin{cases} -\sigma\dot v^1 - \dot v^2 = 0 & \left(\cdot = \frac{d}{ds}\right) \\ -\sigma\dot v^2 - p(v^1)^\cdot = \ddot v^2. \end{cases}$$
An integration using (20) gives
$$(22) \qquad \begin{cases} \sigma v^1 + v^2 = \sigma v^1_l + v^2_l = \sigma v^1_r + v^2_r \\ \dot v^2 = \sigma(v^2_l-v^2) + p(v^1_l) - p(v^1) = \sigma(v^2_r-v^2) + p(v^1_r) - p(v^1), \end{cases}$$
for $u_l=(v^1_l,v^2_l)$, $u_r=(v^1_r,v^2_r)$. In particular,
$$\begin{cases} \sigma v^1_l + v^2_l = \sigma v^1_r + v^2_r \\ \sigma v^2_l + p(v^1_l) = \sigma v^2_r + p(v^1_r). \end{cases}$$
Solving these equations for $\sigma$, we obtain
$$(23) \qquad \sigma^2 = \frac{p(v^1_r)-p(v^1_l)}{v^1_r - v^1_l}.$$
Suppose hereafter $v^1_r > v^1_l$. In view of (18) we can take $\sigma>0$. In this situation the Liu entropy criterion reads
$$(24) \qquad \frac{p(z_1)-p(v^1_l)}{z_1-v^1_l} > \frac{p(v^1_r)-p(v^1_l)}{v^1_r-v^1_l}$$
for all $z$ on the curve $S_k(u_l)$ between $u_l$ and $u_r$, $z=(z_1,z_2)$.

We now claim the system of ODE (22), with asymptotic boundary conditions (20), has a solution if and only if the entropy condition (24) holds. To confirm this, combine the two equations in (22) to eliminate $v^2$:
$$\dot v^1 = \frac{p(v^1)-p(v^1_l)}{\sigma} - \sigma(v^1-v^1_l) =: g(v^1).$$
Now $g(v^1_l)=0$ and $g(v^1_r)=0$, according to (23). Thus in order that the ODE (24) have a solution, with $\lim_{s\to-\infty} v^1 = v^1_l$, $\lim_{s\to\infty} v^1 = v^1_r$, we require
$$g(z_1) > 0 \quad \text{for } v^1_l < z_1 < v^1_r.$$
But this is precisely the entropy criterion (24). A similar calculation works if $v^1_r < v^1_l$. $\square$
\end{example}

\subsection{Entropy/entropy-flux pairs}\label{sec:entropy-entropy-flux-pairs}

Both Lax's and Liu's entropy criteria provide restrictions on possible left and right hand states joined by a shock wave (or a traveling wave for the viscous approximation). It is however of considerable interest to widen still further the entropy criteria, so as to apply to more general integral solutions of our conservation laws.

One idea is to require an integral solution satisfy certain ``entropy-type'' inequalities.

\begin{definition}[Entropy/entropy-flux pair, system case]
\label{def:entropy-entropy-flux-pair-system}
\dcref{ch11:11.4-def-2}
\group{Conservation Laws}
Two smooth functions $\Phi, \Psi:\mathbb{R}^m\to\mathbb{R}$ comprise an \textit{entropy/entropy-flux pair} for the conservation law $u_t + \mathbf{F}(u)_x = 0$ provided

$$(25) \qquad \Phi \text{ is convex}$$

and

$$(26) \qquad D\Phi(z)D\mathbf{F}(z) = D\Psi(z) \quad (z\in\mathbb{R}^m).$$
\end{definition}

To motivate condition (26), suppose for the moment $\mathbf{u}$ is a smooth solution of the system of PDE $u_t + \mathbf{F}(u)_x = 0$. We then compute

$$(27) \qquad \Phi(\mathbf{u})_t + \Psi(\mathbf{u})_x = D\Phi(\mathbf{u})\cdot\mathbf{u}_t + D\Psi(\mathbf{u})\cdot\mathbf{u}_x = (-D\Phi(\mathbf{u})D\mathbf{F}(\mathbf{u}) + D\Psi(\mathbf{u}))\cdot\mathbf{u}_x = 0$$

by (26). This computation says the quantity $\Phi(\mathbf{u})$ satisfies a scalar conservation law, with flux $\Psi(\mathbf{u})$.

Now in general integral solutions of (1) will not be smooth enough, owing to shocks and other irregularities, to justify the foregoing computation. The new idea is instead to replace (27) with an \textit{inequality}:

$$(28) \qquad \Phi(\mathbf{u})_t + \Psi(\mathbf{u})_x \leq 0 \quad \text{in } \mathbb{R}\times(0,\infty).$$

In applications $\Phi(\mathbf{u})$ will sometimes be the negative of physical entropy and $\Psi(\mathbf{u})$ the entropy flux. The inequality (28) therefore asserts entropy evolves according to its flux, but may also undergo sharp increases, for instance along shocks.

Let us hereafter rigorously understand (28) to \textit{mean}

$$(29) \qquad \int_0^\infty\int_{-\infty}^\infty [\Phi(\mathbf{u})v_t + \Psi(\mathbf{u})v_x]\,dx\,dt \geq 0 \quad \text{for each } v\in C_c^\infty(\mathbb{R}\times(0,\infty)), \ v\geq 0.$$

We consider once more the initial-value problem

$$(30) \qquad \begin{cases} u_t + \mathbf{F}(u)_x = 0 & \text{in } \mathbb{R}\times(0,\infty) \\ u = g & \text{on } \mathbb{R}\times\{t=0\}. \end{cases}$$

\begin{definition}[Entropy solution, system case]
\label{def:entropy-solution-system}
\dcref{ch11:11.4-def-3}
\group{Conservation Laws}
\uses{def:integral-solution-system,def:entropy-entropy-flux-pair-system}
We call $u$ an \textit{entropy solution} of (30) provided $u$ is an integral solution and $u$ satisfies the inequalities (29) for each entropy/entropy-flux pair $(\Phi,\Psi)$.
\end{definition}

Let us now attempt to build for general initial data $\mathbf{g}$ an entropy solution. As in \S11.4.1 we expect such a ``physically correct'' solution $\mathbf{u}$ to be a limit of solutions $\mathbf{u}^\varepsilon$ of approximating viscous problems

$$(31) \qquad \begin{cases} u^\varepsilon_t + F(u^\varepsilon)_x - \varepsilon u^\varepsilon_{xx} = 0 & \text{in } \mathbb{R}\times(0,\infty) \\ u^\varepsilon = g & \text{on } \mathbb{R}\times\{t=0\}. \end{cases}$$

We assume $u^\varepsilon$ is a smooth solution of (31), converging to 0 as $|x|\to\infty$ sufficiently rapidly to justify the calculations below. Let us further suppose $\{u^\varepsilon\}_{0<\varepsilon<1}$ is uniformly bounded in $L^\infty$ and furthermore

$$(32) \qquad u^\varepsilon \to u \quad \text{a.e. as } \varepsilon\to 0$$

for some limit function $\mathbf{u}$. (In practice it is extremely difficult to verify this a.e. convergence.)

\begin{theorem}[Entropy and vanishing viscosity]
\label{thm:entropy-vanishing-viscosity}
\dcref{ch11:11.4-thm-2}
\group{Conservation Laws}
\uses{def:entropy-solution-system,def:entropy-entropy-flux-pair-system}
The function $u$ is an entropy solution of the conservation law (30).
\end{theorem}

\begin{proof}
1. Choose any smooth entropy/entropy-flux pair $(\Phi,\Psi)$. Left multiplying (30) by $D\Phi(u^\varepsilon)$ and recalling (26), we compute

$$(33) \qquad \Phi(u^\varepsilon)_t + \Psi(u^\varepsilon)_x = \varepsilon D\Phi(u^\varepsilon)u^\varepsilon_{xx} = \varepsilon\Phi(u^\varepsilon)_{xx} - \varepsilon(D^2\Phi(u^\varepsilon)u^\varepsilon_x)\cdot u^\varepsilon_x.$$

As $\Phi$ is convex,

$$(34) \qquad (D^2\Phi(u^\varepsilon)u^\varepsilon_x)\cdot u^\varepsilon_x \geq 0.$$

2. Multiply (33) by $v\in C_c^\infty(\mathbb{R}\times(0,\infty))$, $v\geq 0$. We integrate by parts and discover:

$$\int_0^\infty\int_{-\infty}^\infty \Phi(u^\varepsilon)v_t + \Psi(u^\varepsilon)v_x\,dx\,dt = \int_0^\infty\int_{-\infty}^\infty \varepsilon(D^2\Phi(u^\varepsilon)u^\varepsilon_x)\cdot u^\varepsilon_x v - \varepsilon\Phi(u^\varepsilon)v_{xx}\,dx\,dt \geq -\int_0^\infty\int_{-\infty}^\infty \varepsilon\Phi(u^\varepsilon)v_{xx}\,dx\,dt,$$

the last inequality holding in view of (34) and the nonnegativity of $v$.

Now let $\varepsilon\to 0$. Recalling (32) and the Dominated Convergence Theorem, we obtain
$$\int_0^\infty\int_{-\infty}^\infty \Phi(u)v_t + \Psi(u)v_x\,dx\,dt \geq 0.$$
Thus $\mathbf{u}$ verifies the entropy/entropy-flux inequalities (29). If $\Phi$ and $\Psi$ are not smooth, we obtain the same conclusion after an approximation.

3. Finally fix $v\in C_c^\infty(\mathbb{R}\times[0,\infty);\mathbb{R}^m)$ and take the dot product of the PDE in (31) with $v$. After integrating by parts we obtain
$$\int_0^\infty\int_{-\infty}^\infty \mathbf{u}^\varepsilon\cdot v_t + F(\mathbf{u}^\varepsilon)v_x + \varepsilon\mathbf{u}^\varepsilon\cdot v_{xx}\,dx\,dt + \int_{-\infty}^\infty \mathbf{g}\cdot v\,dx\Big|_{t=0} = 0.$$
We send $\varepsilon\to 0$, to deduce $\mathbf{u}$ is an integral solution of (30). $\square$
\end{proof}

\begin{example}[Entropy pairs for a scalar conservation law]
\label{ex:scalar-entropy-flux-existence}
\dcref{ch11:11.4-ex-2}
\group{Conservation Laws}
\uses{def:entropy-entropy-flux-pair-system}
In the case of a scalar conservation law (i.e. $m=1$), for any convex $\Phi$ we can find a corresponding flux function $\Psi$, namely
$$\Psi(z) = \int_{z_0}^z \Phi'(w)F'(w)\,dw \quad (z\in\mathbb{R}).$$
See \S11.4.3 following for an application. $\square$
\end{example}

\begin{example}[Entropy pair for the $p$-system]
\label{ex:p-system-entropy-pair}
\dcref{ch11:11.4-ex-3}
\group{Conservation Laws}
\uses{def:entropy-entropy-flux-pair-system,ex:p-system-definition}
For the $p$-system we have $m=2$. To verify (25), (26) we must find $\Phi,\Psi$, with $\Phi$ convex and
$$(\Phi_{z_1},\Phi_{z_2})\begin{pmatrix} 0 & -1 \\ -p'(z_1) & 0 \end{pmatrix} = \begin{pmatrix} \Psi_{z_1} \\ \Psi_{z_2} \end{pmatrix}.$$
A solution is
$$\Phi(z) = \frac{z_2^2}{2} + \int_0^{z_1} p(w)\,dw, \quad \Psi(z) = -p(z_1)z_2 \quad (z\in\mathbb{R}^2).$$
Note $\Phi$ is convex, since $p'>0$. $\square$
\end{example}

See Problems 4, 5 for other examples.

\subsection{Uniqueness for a scalar conservation law}\label{sec:uniqueness-scalar-conservation-law-entropy}

As a further illustration of the ideas in \S11.4.2, let us now consider again the initial-value problem for a \textit{scalar} conservation law

$$(35) \qquad \begin{cases} u_t + F(u)_x = 0 & \text{in } \mathbb{R}\times(0,\infty) \\ u = g & \text{on } \mathbb{R}\times\{t=0\}. \end{cases}$$

Hence the unknown $u=u(x,t)$ is real-valued and $F:\mathbb{R}\to\mathbb{R}$ is a given smooth flux function.

In \S3.4 we carefully studied the problem (35), making use of the primary assumption that $F$ be strictly convex, to derive the Lax--Oleinik formula (\cref{thm:lax-oleinik-formula}; see \S3.4.2). Let us now drop the assumption that $F$ be convex and devise an appropriate notion of weak solution. As above, we introduce entropies:

\begin{definition}[Entropy/entropy-flux pair, scalar case]
\label{def:entropy-entropy-flux-pair-scalar}
\dcref{ch11:11.4-def-4}
\group{Conservation Laws}
Two smooth functions $\Phi,\Psi:\mathbb{R}\to\mathbb{R}$ comprise an \textit{entropy/entropy-flux pair} for the conservation law $u_t + F(u)_x = 0$ provided

$$(36) \qquad \Phi \text{ is convex}$$

and

$$(37) \qquad \Phi'(z)F'(z) = \Psi'(z) \quad (z\in\mathbb{R}).$$
\end{definition}

As noted in \cref{ex:scalar-entropy-flux-existence} above, for each convex $\Phi$ there exists a corresponding flux $\Psi$.

The entropy condition for $u$ reads

$$\Phi(u)_t + \Psi(u)_x \leq 0 \quad \text{on } \mathbb{R}\times(0,\infty)$$

for each entropy/entropy-flux pair $\Phi,\Psi$. This means

$$(38) \qquad \left\{\begin{array}{l} \int_0^\infty\int_{-\infty}^\infty \Phi(u)v_t + \Psi(u)v_x\,dx\,dt \geq 0 \\ \text{for each } v\in C_c^\infty(\mathbb{R}\times(0,\infty)), \ v\geq 0. \end{array}\right.$$

\begin{definition}[Entropy solution, scalar case]
\label{def:entropy-solution-scalar-new}
\dcref{ch11:11.4-def-5}
\group{Conservation Laws}
\uses{def:entropy-entropy-flux-pair-scalar}
We call $u\in C([0,\infty), L^1(\mathbb{R}))\cap L^\infty(\mathbb{R}\times(0,\infty))$ an \textit{entropy solution} of (35) provided $u$ satisfies the inequalities (38) for each entropy/entropy-flux pair $(\Phi,\Psi)$, and $u(\cdot,t)\to g$ in $L^1$ as $t\to 0$.
\end{definition}

\begin{remark}[Consistency with the scalar entropy solution of Chapter 3]
\label{rem:entropy-solution-scalar-consistency}
\dcref{ch11:11.4-rem-1}
\group{Conservation Laws}
\uses{def:entropy-solution-scalar-new,def:integral-solution-conservation-law}
(i) This definition supersedes our earlier definition of ``entropy solution'' in \S3.4.3 (\cref{def:entropy-solution-conservation-law}).

(ii) Taking $\Phi(z)=\pm z$, $\Psi(z)=\pm F(z)$ in (38), we deduce
$$\int_0^\infty\int_{-\infty}^\infty uv_t + F(u)v_x\,dx\,dt = 0$$
for all $v\geq 0$ and thus for all $v\in C_c^1(\mathbb{R}\times(0,\infty))$. It is an exercise to prove then that
$$\int_0^\infty\int_{-\infty}^\infty uv_t + F(u)v_x\,dx\,dt + \int_{-\infty}^\infty gv\,dx\Big|_{t=0} = 0$$
for all $v\in C_c^1(\mathbb{R}\times[0,\infty))$, since $u(\cdot,t)\to g$ in $L^1$. \textit{Thus an entropy solution is an integral solution.} $\square$
\end{remark}

We discussed in \S11.4.2 the construction of an entropy solution, and we now prove uniqueness.

\begin{theorem}[Uniqueness of entropy solutions for a single conservation law]
\label{thm:uniqueness-entropy-solutions-scalar-new}
\dcref{ch11:11.4-thm-3}
\group{Conservation Laws}
\uses{def:entropy-solution-scalar-new,thm:uniqueness-viscosity-solution-hj}
There exists---up to a set of measure zero---at most one entropy solution of (35).
\end{theorem}

As in the proof of \cref{thm:uniqueness-viscosity-solution-hj} in \S10.2, the basic idea will be to ``double the variables'' in the problem.

\begin{proof}
1. Let $u$ be an entropy solution of (35). Then

$$(39) \qquad \int_0^\infty\int_{-\infty}^\infty \Phi(u)v_t + \Psi(u)v_x\,dx\,dt \geq 0$$

for all $v\in C_c^\infty(\mathbb{R}\times(0,\infty))$, $v\geq 0$, where $\Phi$ is smooth, convex and

$$\Psi(z) = \int_{z_0}^z \Phi'(w)F'(w)\,dw$$

for any $z_0$. Fix $\alpha\in\mathbb{R}$ and take

$$(40) \qquad \Phi_k(z) := \beta_k(z-\alpha) \quad (z\in\mathbb{R}),$$

where for each $k=1,\ldots$, $\beta_k:\mathbb{R}\to\mathbb{R}$ is smooth, convex and

$$\begin{cases} \beta_k(z) \to |z| & \text{uniformly} \\ \beta_k'(z) \to \operatorname{sgn}(z) & \text{boundedly, a.e.} \end{cases}$$

Thus $\Phi_k(z)\to|z-\alpha|$ uniformly for $z\in\mathbb{R}$. A flux corresponding to (40) is

$$\Psi_k(z) = \int_\alpha^z \beta_k(w-\alpha)F'(w)\,dw.$$

Consequently for each $z$

$$\Psi_k(z) \to \int_\alpha^z \operatorname{sgn}(w-\alpha)F'(w)\,dw = \operatorname{sgn}(z-\alpha)(F(z)-F(\alpha)).$$

Putting $\Phi_k$, $\Psi_k$ into (39) and sending $k\to\infty$, we deduce

$$(41) \qquad \int_0^\infty\int_{-\infty}^\infty |u-\alpha|v_t + \operatorname{sgn}(u-\alpha)(F(u)-F(\alpha))v_x\,dx\,dt \geq 0$$

for each $\alpha\in\mathbb{R}$ and $v$ as above.

2. Next let $\tilde u$ be another entropy solution. Then

$$(42) \qquad \int_0^\infty\int_{-\infty}^\infty |\tilde u-\bar a||\bar v_s + \operatorname{sgn}(\tilde u-\bar a)(F(\tilde u)-F(\bar a))\bar v_y\,dy\,ds \geq 0$$

where $\bar a\in\mathbb{R}$ and $\bar v\in C_c^\infty(\mathbb{R}\times(0,\infty))$, $\bar v\geq 0$.

Now let $w\in C_c^\infty(\mathbb{R}\times\mathbb{R}\times(0,\infty)\times(0,\infty))$, $w\geq 0$, $w=w(x,y,t,s)$. Fixing $(y,s)\in\mathbb{R}\times(0,\infty)$, we take $\alpha=\bar u(y,s)$, $v(x,t)=w(x,y,t,s)$ in (41). Integrating with respect to $y,s$, we produce the inequality

$$(43) \qquad \int_0^\infty\int_0^\infty\int_{-\infty}^\infty\int_{-\infty}^\infty |u(x,t)-\bar u(y,s)|w_t + \operatorname{sgn}(u(x,t)-\bar u(y,s))(F(u(x,t))-F(\bar u(y,s)))w_x \, dx\,dy\,dt\,ds \geq 0.$$

Likewise, for each fixed $(x,t)\in\mathbb{R}\times(0,\infty)$ we take $\bar a=u(x,t)$, $\bar v(y,s)=w(x,y,t,s)$ in (42). Integrating with respect to $x,t$ gives:

$$(44) \qquad \int_0^\infty\int_0^\infty\int_{-\infty}^\infty\int_{-\infty}^\infty |\bar u(y,s)-u(x,t)|w_s + \operatorname{sgn}(\bar u(y,s)-u(x,t))(F(\bar u(y,s))-F(u(x,t)))w_y \, dx\,dy\,dt\,ds \geq 0.$$

Add (43), (44):

$$(45) \qquad \int_0^\infty\int_0^\infty\int_{-\infty}^\infty\int_{-\infty}^\infty |u(x,t)-\bar u(y,s)|(w_t+w_s) + \operatorname{sgn}(u(x,t)-\bar u(y,s))(F(u(x,t))-F(\bar u(y,s)))(w_x+w_y) \, dx\,dy\,dt\,ds \geq 0.$$

3. We design as follows a clever choice for $w$ in (45). Select $\eta$ to be a standard mollifier as in \S C.4 (with $n=1$) and, as usual, write $\eta_\epsilon(x) = \frac{1}{\epsilon}\eta\left(\frac{x}{\epsilon}\right)$. Take

$$w(x,y,t,s) := \eta_\epsilon\left(\frac{x-y}{2}\right)\eta_\epsilon\left(\frac{t-s}{2}\right)\phi\left(\frac{x+y}{2},\frac{t+s}{2}\right),$$

where $\phi\in C_c^\infty(\mathbb{R}\times(0,\infty))$, $\phi\geq 0$. We insert this choice of $w$ into (45) and thereby obtain:

$$(46) \qquad \int_0^\infty\int_0^\infty\int_{-\infty}^\infty\int_{-\infty}^\infty \left\{|u(x,t)-\bar u(y,s)|\phi_t\left(\frac{x+y}{2},\frac{t+s}{2}\right) + \operatorname{sgn}(u(x,t)-\bar u(y,s))(F(u(x,t))-F(\bar u(y,s)))\phi_x\left(\frac{x+y}{2},\frac{t+s}{2}\right)\right\} \eta_\epsilon\left(\frac{x-y}{2}\right)\eta_\epsilon\left(\frac{t-s}{2}\right)dx\,dy\,dt\,ds \geq 0.$$

Change variables by writing

$$\bar x = \frac{x+y}{2}, \quad \bar t = \frac{t+s}{2}, \qquad \bar y = \frac{x-y}{2}, \quad \bar s = \frac{t-s}{2}.$$

Then (46) implies

$$(47) \qquad \int_{-\infty}^\infty\int_{-\infty}^\infty G(\bar y,\bar s)\eta_\epsilon(\bar y)\eta_\epsilon(\bar s)\,d\bar y\,d\bar s \geq 0,$$

where

$$(48) \qquad G(\bar y,\bar s) := \int_0^\infty\int_{-\infty}^\infty |u(\bar x+\bar y,\bar t+\bar s) - \bar u(\bar x-\bar y,\bar t-\bar s)|\phi_t(\bar x,\bar t) + \operatorname{sgn}(u(\bar x+\bar y,\bar t+\bar s)-\bar u(\bar x-\bar y,\bar t-\bar s))(F(u(\bar x+\bar y,\bar t+\bar s))-F(\bar u(\bar x-\bar y,\bar t-\bar s)))\phi_x(\bar x,\bar t)\,d\bar x\,d\bar t.$$

Now $u(\bar x+\bar y,\bar t+\bar s) \to u(\bar x,\bar t)$, $\bar u(\bar x-\bar y,\bar t-\bar s) \to \bar u(\bar x,\bar t)$ in $L^1_{\text{loc}}$ as $\bar y,\bar s\to 0$. Since the mappings $(a,b)\mapsto |a-b|$, $\operatorname{sgn}(a-b)(F(a)-F(b))$ are Lipschitz continuous, we deduce upon letting $\epsilon\to 0$ in (47) that

$$\int_0^\infty\int_{-\infty}^\infty |u(\bar x,\bar t)-\bar u(\bar x,\bar t)|\phi_t(\bar x,\bar t) + \operatorname{sgn}(u(\bar x,\bar t)-\bar u(\bar x,\bar t))(F(u(\bar x,\bar t))-F(\bar u(\bar x,\bar t)))\phi_x(\bar x,\bar t)\,d\bar x\,d\bar t \geq 0.$$

Rewriting $x=\bar x$, $t=\bar t$, we have therefore

$$(49) \qquad \int_0^\infty\int_{-\infty}^\infty a(x,t)\phi_t(x,t) + b(x,t)\phi_x(x,t)\,dx\,dt \geq 0,$$

for
$$\begin{cases} a(x,t) := |u(x,t)-\bar u(x,t)| \\ b(x,t) := \operatorname{sgn}(u(x,t)-\bar u(x,t))(F(u(x,t))-F(\bar u(x,t))). \end{cases}$$

4. We now employ the inequality (49) to establish the $L^1$-contraction inequalities

$$(50) \qquad \begin{cases} \int_{-\infty}^\infty |u(x,t)-\bar u(x,t)|\,dx \leq \int_{-\infty}^\infty |u(x,s)-\bar u(x,s)|\,dx \\ \text{for a.e. } 0\leq s\leq t. \end{cases}$$

To prove this assertion, we take $0<s<t$, $r>0$, and let $\phi(x,t)=\alpha(x)\beta(t)$ in (49), where

$$\begin{cases} \alpha:\mathbb{R}\to\mathbb{R} \text{ is smooth,} \\ \alpha(x)=1 \text{ if } |x|\leq r, \quad \alpha(x)=0 \text{ if } |x|\geq r+1, \\ |\alpha'(x)|\leq 2 \end{cases}$$

and

$$\begin{cases} \beta:\mathbb{R}\to\mathbb{R} \text{ is Lipschitz,} \\ \beta(\tau)=0 \text{ if } 0\leq\tau\leq s \text{ or } \tau\geq t+\delta, \\ \beta(\tau)=1 \text{ if } s+\delta\leq\tau\leq t, \\ \beta \text{ is linear on } [s,s+\delta] \text{ and } [t,t+\delta], \end{cases}$$

for $0<\delta<t-s$. We deduce

$$\frac{1}{\delta}\int_s^{s+\delta}\int_{-\infty}^\infty a(x,\tau)\alpha(x)\,dx\,d\tau \geq \frac{1}{\delta}\int_t^{t+\delta}\int_{-\infty}^\infty a(x,\tau)\alpha(x)\,dx\,d\tau - \int_s^{t+\delta}\int_{\{r\leq|x|\leq r+1\}} b(x,\tau)\alpha'(x)\beta(\tau)\,dx\,d\tau.$$

Let $r\to\infty$:

$$\frac{1}{\delta}\int_t^{t+\delta}\int_{-\infty}^\infty a(x,\tau)\,dx\,d\tau \leq \frac{1}{\delta}\int_s^{s+\delta}\int_{-\infty}^\infty a(x,\tau)\,dx\,d\tau.$$

Next let $\delta\to 0$ to deduce (50) for a.e. $0\leq s\leq t$.

5. In light of (50) and the fact $u(\cdot,t), \bar u(\cdot,t) \to g$ in $L^1$ as $t\to 0$, we at last conclude $u=\bar u$ a.e. $\square$
\end{proof}

\section{Problems}\label{sec:chapter11-problems}

%ORIGINAL-NUMBER: 1
\begin{exercise}
Verify that the shallow water equations (\cref{ex:shallow-water-equations}) form a strictly hyperbolic system, provided $\varphi > 0$.
\end{exercise}

%ORIGINAL-NUMBER: 2
\begin{exercise}
Define for $z\in\mathbb{R}$, $z\neq 0$, the matrix function
$$\mathbf{B}(z) := e^{-1/z^2}\begin{pmatrix} \cos(2/z) & \sin(2/z) \\ \sin(2/z) & -\cos(2/z) \end{pmatrix},$$
and set $\mathbf{B}(0)=0$. Show that $\mathbf{B}$ is $C^\infty$ and has real eigenvalues, but we cannot find unit-length right eigenvectors $\{\mathbf{r}_1(z),\mathbf{r}_2(z)\}$ depending continuously on $z$ near $0$. What happens to the eigenspaces as $z\to 0$?
\end{exercise}

%ORIGINAL-NUMBER: 3
\begin{exercise}
Confirm that the functions $w^1, w^2$ computed for the barotropic gas dynamics in \S11.3.1 are indeed Riemann invariants.
\end{exercise}

%ORIGINAL-NUMBER: 4
\begin{exercise}
Suppose that $\Phi$ is an entropy for the shallow water equations. Prove
$$\frac{\partial^2\Phi}{\partial v^2} = \varphi\frac{\partial^2\Phi}{\partial\rho^2}.$$
\end{exercise}

%ORIGINAL-NUMBER: 5
\begin{exercise}
Show that $\Phi = \rho v^2/2 + P(\rho)$ is an entropy for the barotropic Euler equations (from \S11.3.1), provided $P''(\rho) = p'(\rho)/\rho$, $\rho>0$. Confirm that $\Phi$ is convex in the proper variables. What is the corresponding entropy flux $\Psi$?
\end{exercise}

%ORIGINAL-NUMBER: 6
\begin{exercise}
Assume that $u$ is an entropy solution of the scalar conservation law $u_t + F(u)_x = 0$, and that, as in \S3.4.1, $u$ is smooth on either side of a curve $\{x=s(t)\}$.

(i) Prove that along this curve the left and right hand limits of $u$ satisfy the relations:
$$F(z) \geq \frac{F(u_r)-F(u_l)}{u_r-u_l}(z-u_r) + F(u_r) \quad \text{if } u_l \leq z \leq u_r,$$
and
$$F(z) \leq \frac{F(u_r)-F(u_l)}{u_r-u_l}(z-u_r) + F(u_r) \quad \text{if } u_r \leq z \leq u_l.$$
This is the \textit{condition E}.

(ii) What does condition E imply if $F$ is uniformly convex?
\end{exercise}

\section{References}\label{sec:chapter11-references}

T.-P. Liu wrote the initial draft of this entire chapter.

\textbf{Section 11.1.} The shallow water equations are discussed in LeVeque \cite{LeVeque1992}. P. Colella helped me with the calculations in \S11.1.2. The proof of \cref{thm:smooth-dependence-eigenvalues-eigenvectors} follows suggestions of W. Han. See Courant--Friedrichs \cite{CourantFriedrichs1976}, Xiao--Zhang \cite{XiaoZhang1989}, Zheng \cite{Zheng2001} for more.

\textbf{Section 11.3.} I followed Logan \cite{Logan1987} for the example.

\textbf{Section 11.4.} See Majda--Pego \cite{MajdaPego1985}, Conlon \cite{Conlon1980}, Smoller \cite{Smoller1983} for more about viscous traveling waves. The uniqueness theorem in \S11.4.3 is due to Kru\v{z}kov.

\textbf{Section 11.5.} Problem 2 is taken from Kato \cite{Kato1980}, p. 111.
